\documentclass[openany]{book}

\usepackage[margin=1in]{geometry}
\usepackage{amsmath,amsfonts,amsthm, amssymb}
\usepackage{yhmath}
\usepackage{mathrsfs}
\usepackage{mathtools}
\usepackage{xcolor}
\usepackage{graphicx}
\usepackage{comment}
\usepackage{tikz-cd}
\usepackage{quiver}
\renewcommand{\familydefault}{ppl}
\newcommand{\R}{\mathbb{R}}
\newcommand{\E}{\mathbb{E}}
\newcommand{\Z}{\mathbb{Z}}
\newcommand{\CC}{\mathbb{C}}
\newcommand{\F}{\mathbb{F}}
\newcommand{\la}{\langle}
\newcommand{\ra}{\rangle}
\newcommand{\colim}{\text{colim}}
\DeclareMathOperator{\im}{im}
\let\oldemptyset\emptyset
\let\emptyset\varnothing
\newcommand{\tor}{\text{Tor}}
\newcommand{\id}{\text{id}}
\newcommand{\ext}{\text{Ext}}
\newcommand{\ptop}{\text{PTop}}
\newcommand{\pt}{\text{pt}}
\newcommand{\ach}{\text{Ach}}
\newcommand{\Q}{\mathbb{Q}}


\usepackage{thmtools,thm-restate}

% Fixing mdframed skip below
% See https://tex.stackexchange.com/a/292090/143086
\usepackage[framemethod=TikZ]{mdframed}
\usepackage{xpatch}
\makeatletter
\xpatchcmd{\endmdframed}
	{\aftergroup\endmdf@trivlist\color@endgroup}
	{\endmdf@trivlist\color@endgroup\@doendpe}
	{}{}
\makeatother

\definecolor{huilightpink}{HTML}{fff2fe}
\definecolor{huidarkpink}{HTML}{d955b7}
\declaretheoremstyle[
	mdframed={
		backgroundcolor=huilightpink,
		linecolor=huidarkpink,
		rightline=false,
		topline=false,
		bottomline=false,
		linewidth=2pt,
		innertopmargin=5pt,
		innerbottommargin=8pt,
		innerleftmargin=8pt,
		leftmargin=-2pt,
		skipbelow=2pt,
		nobreak
	},
	headfont=\normalfont\bfseries\color{huidarkpink}
]{huipinkbox}
\declaretheorem[style=huipinkbox,name=Theorem,within=chapter]{thm}
\declaretheorem[style=huipinkbox,name=Theorem,sibling=thm]{theorem}




\begin{comment}
\definecolor{huilightyellow}{HTML}{fff5d6}
\definecolor{huidarkyellow}{HTML}{fcad03}
\declaretheoremstyle[
	mdframed={
		backgroundcolor=huilightyellow,
		linecolor=huidarkyellow,
		rightline=false,
		topline=false,
		bottomline=false,
		linewidth=2pt,
		innertopmargin=5pt,
		innerbottommargin=8pt,
		innerleftmargin=8pt,
		leftmargin=-2pt,
		skipbelow=2pt,
		nobreak
	},
	headfont=\normalfont\bfseries\color{huidarkyellow}
]{huiyellowbox}
\declaretheorem[style=huiyellowbox,name=Proposition,within=chapter]{prop}
\end{comment}



\definecolor{huilightpurple}{HTML}{faf2ff}
\definecolor{huidarkpurple}{HTML}{912ed9}
\declaretheoremstyle[
	mdframed={
		backgroundcolor=huilightpurple,
		linecolor=huidarkpurple,
		rightline=false,
		topline=false,
		bottomline=false,
		linewidth=2pt,
		innertopmargin=5pt,
		innerbottommargin=8pt,
		innerleftmargin=8pt,
		leftmargin=-2pt,
		skipbelow=2pt,
		nobreak
	},
	headfont=\normalfont\bfseries\color{huidarkpurple}
]{huipurplebox}
\declaretheorem[style=huipurplebox,name=Proposition,within=chapter]{prop}



% \definecolor{huilightpurple}{HTML}{faf2ff}
% \definecolor{huidarkpurple}{HTML}{912ed9}
% \declaretheoremstyle[
% 	mdframed={
% 		backgroundcolor=huilightpurple,
% 		linecolor=huidarkpurple,
% 		rightline=false,
% 		topline=false,
% 		bottomline=false,
% 		linewidth=2pt,
% 		innertopmargin=5pt,
% 		innerbottommargin=8pt,
% 		innerleftmargin=8pt,
% 		leftmargin=-2pt,
% 		skipbelow=2pt,
% 		nobreak
% 	},
% 	headfont=\normalfont\bfseries\color{huidarkpurple}
% ]{huipurplebox}
\declaretheorem[style=huipurplebox,name=Lemma,within=chapter]{lem}


\definecolor{lightpink}{HTML}{f0f6fc}
\definecolor{darkpink}{HTML}{2c72b8}
\declaretheoremstyle[
	mdframed={
		backgroundcolor=lightpink,
		linecolor=darkpink,
		rightline=false,
		topline=false,
		bottomline=false,
		linewidth=2pt,
		innertopmargin=5pt,
		innerbottommargin=8pt,
		innerleftmargin=8pt,
		leftmargin=-2pt,
		skipbelow=2pt,
		nobreak
	},
	headfont=\normalfont\bfseries\color{darkpink}
]{pinkbox}
\declaretheorem[style=pinkbox,name=Definition,within=chapter]{defn}


\definecolor{huilightblue}{HTML}{edf9ff}
\definecolor{huidarkblue}{HTML}{4b79db}
\declaretheoremstyle[
	mdframed={
		backgroundcolor=huilightblue,
		linecolor=huidarkblue,
		rightline=false,
		topline=false,
		bottomline=false,
		linewidth=2pt,
		innertopmargin=5pt,
		innerbottommargin=8pt,
		innerleftmargin=8pt,
		leftmargin=-2pt,
		skipbelow=2pt,
		nobreak
	},
	headfont=\normalfont\bfseries\color{huidarkblue}
]{huiblueblox}
\declaretheorem[style=huiblueblox,name=Example,within=chapter]{example}



% \definecolor{huilightblue}{HTML}{edf9ff}
% \definecolor{huidarkblue}{HTML}{4b79db}
% \declaretheoremstyle[
% 	mdframed={
% 		backgroundcolor=huilightblue,
% 		linecolor=huidarkblue,
% 		rightline=false,
% 		topline=false,
% 		bottomline=false,
% 		linewidth=2pt,
% 		innertopmargin=5pt,
% 		innerbottommargin=8pt,
% 		innerleftmargin=8pt,
% 		leftmargin=-2pt,
% 		skipbelow=2pt,
% 		nobreak
% 	},
% 	headfont=\normalfont\bfseries\color{huidarkblue}
% ]{huiblueblox}
% \declaretheorem[style=huiblueblox,name=Example,within=chapter]{example}

% \declaretheoremstyle[
% 	mdframed={
% 		backgroundcolor=huilightblue,
% 		linecolor=huidarkblue,
% 		rightline=false,
% 		topline=false,
% 		bottomline=false,
% 		linewidth=2pt,
% 		innertopmargin=5pt,
% 		innerbottommargin=8pt,
% 		innerleftmargin=8pt,
% 		leftmargin=-2pt,
% 		skipbelow=2pt,
% 		nobreak
% 	},
% 	headfont=\normalfont\bfseries\color{huidarkblue}
% ]{huiblueblox}
\declaretheorem[style=huiblueblox,name=Problem,within=chapter]{prob}



% \declaretheoremstyle[
% 	mdframed={
% 		backgroundcolor=huilightblue,
% 		linecolor=huidarkblue,
% 		rightline=false,
% 		topline=false,
% 		bottomline=false,
% 		linewidth=2pt,
% 		innertopmargin=5pt,
% 		innerbottommargin=8pt,
% 		innerleftmargin=8pt,
% 		leftmargin=-2pt,
% 		skipbelow=2pt,
% 		nobreak
% 	},
% 	headfont=\normalfont\bfseries\color{huidarkblue}
% ]{huiblueblox}
\declaretheorem[style=huiblueblox,name=Exercise,within=chapter]{exer}
\declaretheorem[style=huipinkbox, name=Corollary, within=chapter]{cor}





















\newcommand{\nirwarnsymbol}{%
	\begin{tikzpicture}[baseline=(x.base)]
		\draw[rounded corners=.01em] (-.05em,-1.07em)rectangle(.05em,.78em);
		\draw[fill=white,rounded corners=1.3] (0,.75em)--(.75em,0)--(0,-.75em)--(-.75em,0)--cycle;
		\draw[line width=0.2mm, line cap=round](-.4em,-1.07em)--(.4em,-1.07em);
		\node(x) at (0,0em) {};
		% Thank you https://tex.stackexchange.com/a/262510
		\draw[
			line cap=but,
			line join=round,
			x=.5em,
			line width=0.5mm,
			y=1*(height("Z")-\pgflinewidth)*(1-sin(10)),
			rotate=-10,
			rounded corners=1.5pt,
		](-0.57, 0.57) -- (0.57, 0.57) -- (-0.57, -0.57) -- (0.57, -0.57);
	\end{tikzpicture}%
}

%%%%%%%%%%%%%%%%%%%%%%%%%%%%%%%%%%%%%%%%%%%% MARGINS
\usepackage{marginnote}
% Thank you https://tex.stackexchange.com/a/472882
% Makes marginnotes always appear on the left, apparently
%
\makeatletter
\long\def\@mn@@@marginnote[#1]#2[#3]{%
	\begingroup
		\ifmmode\mn@strut\let\@tempa\mn@vadjust\else
			\if@inlabel\leavevmode\fi
			\ifhmode\mn@strut\let\@tempa\mn@vadjust\else\let\@tempa\mn@vlap\fi
		\fi
		\@tempa{%
			\vbox to\z@{%
				\vss
				\@mn@margintest
				\if@reversemargin\if@tempswa
						\@tempswafalse
					\else
						\@tempswatrue
				\fi\fi

					\llap{%
						\vbox to\z@{\kern\marginnotevadjust\kern #3
							\vbox to\z@{%
								\hsize\marginparwidth
								\linewidth\hsize
								\kern-\parskip
								%\mn@parboxrestore
								\marginfont\raggedleftmarginnote\strut\hspace{\z@}%
								\ignorespaces#1\endgraf
								\vss
							}%
							\vss
						}%
						\if@mn@verbose
							\PackageInfo{marginnote}{xpos seems to be \@mn@currxpos}%
						\fi
						\begingroup
							\ifx\@mn@currxpos\relax\else\ifx\@mn@currpos\@empty\else
									\kern\@mn@currxpos
							\fi\fi
							\ifx\@mn@currpage\relax
								\let\@mn@currpage\@ne
							\fi
							\if@twoside\ifodd\@mn@currpage\relax
									\kern-\oddsidemargin
								\else
									\kern-\evensidemargin
								\fi
							\else
								\kern-\oddsidemargin
							\fi
							\kern-1in
						\endgroup
						\kern\marginparsep
					}%
			}%
		}%
	\endgroup
}
\makeatother
%
% Mostly for todonotes
\renewcommand{\marginpar}{\marginnote}
%%%%%%%%%%%%%%%%%%%%%%%%%%%%%%%%%%%%%%%%%%%% /MARGINS

\definecolor{nirlightred}{RGB}{250, 220, 220}
\definecolor{nirdarkred}{HTML}{f40000}
\declaretheoremstyle[
	mdframed={
		backgroundcolor=nirlightred,
		linecolor=nirdarkred,
		rightline=false,
		topline=false,
		bottomline=false,
		linewidth=2pt,
		innertopmargin=5pt,
		innerbottommargin=8pt,
		innerleftmargin=8pt,
		leftmargin=-2pt,
		skipbelow=2pt,
		nobreak
	},
	headfont=\normalfont\bfseries\color{nirdarkred}
]{nirredbox}

% \makeatletter
% \declaretheorem[
% 	style=nirredbox,
% 	name=Warning,
% 	sibling=thm,
% 	% without \leavevmode, the first item in a list gets misformatted
% 	postheadhook={\leavevmode\marginnote{\nirwarnsymbol}[-3pt]%
% 	\ifthmt@thisistheone% restatable makes alignment weird
% 		\hspace{-2.2pt}%
% 	\fi}
% ]{warn}
% \makeatother

\newcommand{\nirideasymbol}{%
	\begin{tikzpicture}[baseline=(x.base)]
		\draw[rounded corners=.01em] (-.05em,-1.07em)rectangle(.05em,.78em);
		\draw[fill=white,rounded corners=1.3] (0,.75em)--(.75em,0)--(0,-.75em)--(-.75em,0)--cycle;
		\draw[line width=0.2mm, line cap=round](-.4em,-1.07em)--(.4em,-1.07em);
		\node(x) at (0,0em) {};
		\node at (0,0em) {{\textbf{!}}};
	\end{tikzpicture}%
}
\renewcommand{\nirwarnsymbol}{%
	\begin{tikzpicture}[baseline=(x.base)]
		\draw[rounded corners=.01em] (-.05em,-1.07em)rectangle(.05em,.78em);
		\draw[fill=white,rounded corners=1.3] (0,.75em)--(.75em,0)--(0,-.75em)--(-.75em,0)--cycle;
		\draw[line width=0.2mm, line cap=round](-.4em,-1.07em)--(.4em,-1.07em);
		\node(x) at (0,0em) {};
		% Thank you https://tex.stackexchange.com/a/262510
		\draw[
			line cap=but,
			line join=round,
			x=.5em,
			line width=0.5mm,
			y=1*(height("Z")-\pgflinewidth)*(1-sin(10)),
			rotate=-10,
			rounded corners=1.5pt,
		](-0.57, 0.57) -- (0.57, 0.57) -- (-0.57, -0.57) -- (0.57, -0.57);
	\end{tikzpicture}%
}
\makeatletter
\declaretheorem[
	style=nirredbox,
	name=Idea,
	sibling=thm,
	% without \leavevmode, the first item in a list gets misformatted
	postheadhook={\leavevmode\marginnote{\nirideasymbol}[-3pt]%
	\ifthmt@thisistheone% restatable makes alignment weird
		\hspace{-2.2pt}%
	\fi}
]{idea}

\declaretheorem[
	style=nirredbox,
	name=Warning,
	sibling=thm,
	% without \leavevmode, the first item in a list gets misformatted
	postheadhook={\leavevmode\marginnote{\nirwarnsymbol}[-3pt]%
	\ifthmt@thisistheone% restatable makes alignment weird
		\hspace{-2.2pt}%
	\fi}
]{warn}
\makeatother

\title{Algebra Qualifying Exam}

\date{\today}
\author{Naruki Masuda, transcribed by Hui Sun}


\begin{document}

\maketitle

\tableofcontents



\newpage

% \chapter{1}

% We will first talk about algebraic sets and spaces.

% \begin{thm}[1.5 Hilbert's Nullstellnsatz, p380]
%     Let $a\subset k[x]$ be an ideal, where $x=(x_1,\dots,x_n)$, let $f\in k[x]$ be a polynomial such that $f(c)=0$ for every zero $c=(c_1,\dots, c_n)$ of $a\subset k[x]$, then there exists $m\in\mathbb{N}$ such that 
%     \begin{equation*}
%         f^m\in a
%     \end{equation*}
%     Note that $f^m\in a \iff f\in\sqrt{a}\supset a$.
% \end{thm}
% \begin{proof}
%     Suppose $f\neq 0$, and $a$ has a zero. 
% \end{proof}
% The following theorem further characterizes ideals.
% \begin{thm}[1.4]
%     Let $a\subset k[x]=k[x_1,\dots,x_n]$, then either $a=k[x]$ or $a$ has a zero in $k^a$. Now we introduce a new variable $Y$, and consider the ideal $a'$ generated by $a$ and $1-Yf$. By the following theorem, we know that $a'$ must be the entire ring $k[X,Y]$. This implies that there exists polynomials $g_j\in k[X,Y], h_i\in a$ such that 
%     \begin{equation*}
%         1=g_0(1-Yf)+g_1h_1+\dots+g_rh_r
%     \end{equation*}
%     Substitute $Y$ with $f^{-1}$, and multiply both sides by $f^m$, 
%     \begin{equation*}
%         f^m=f^m(g_1h_1+\dots+g_rh_r)\in a
%     \end{equation*}
%     and we are done.
% \end{thm}

% \begin{thm}[1.1 p378]
%     Let $k$ be a field, and $k[x]=k[x_1,\dots, x_n]$ be a finitely generated ring. If $\varphi: k\to L$ is an embedding of an algebraically closed field $L$, then there exists an extension $\bar{\varphi}:k[X]\to L$.
% \end{thm}
% An immediate corollary is as follows:
% \begin{cor}[1.2 p389]
%     Let $k$ be a field and $k[x]$ be finitely generated over $k$. If $k[x]$ is a field, then $k[x]$ is algebraically closed over $k$.
% \end{cor}

% \begin{prop}[3.1 p347]
%     Let $A$ be a subring of $B$, and $B$ is integral over $A$, let $\varphi:A\to L$ where $L$ is a field that is algebraically closed, then $\varphi$ extends to a homomorphism of $B$ to $L$.
% \end{prop}

% Next we proved theorem 1.4 above. I cannot do this.



\chapter{Group Theory}
\textcolor{red}{S2013-Q2, S2016-Q1, F2018-Q2, F2001-Q1, F2013-Q}

\textcolor{blue}{S2005-Q1, F2009-Q1}


\section{Sylow Theorems}
We first talk bout semidirect products. Let $G$ be any group, and $N,H$ be subgroups of $G$.
\begin{defn}
For $\varphi:H\to\text{Aut}(N)$, define $N\rtimes H$ by 
\begin{itemize}
    \item[(1)] $N\rtimes_\varphi H=N\times H$ as a set.
    \item[(b)] Equipped with the group structure 
    \begin{equation*}
        (n_1,h_1)\cdot (n_2,h_2)=(n_1\varphi(h_1)n_2, h_1h_2)
    \end{equation*}
    The structure $(N\rtimes_\varphi H, \cdot)$ forms a group. 
\end{itemize}
\end{defn}

\begin{example}
    If $N$ is a normal subgroup of $G$, and $N\cap H=\{e\}$, and $\varphi:H\to\text{Aut}(N)$ where 
    \begin{equation*}
        \varphi: h\mapsto (n\mapsto hnh^{-1})
    \end{equation*}
    (acting by conjugation), and $G=NH$. Then 
    \begin{equation*}
        N\rtimes_\varphi H\to G
    \end{equation*}
    where 
    \begin{equation*}
        (n,h)\mapsto nh
    \end{equation*}
    is a bijective homomorphism homomorphism. Hence 
    \begin{equation*}
        G\cong N\rtimes_\varphi H
    \end{equation*}
\end{example}

% \textcolor{red}{what is happening}

Next we present some divisibility results.
\begin{prop}[Lagrange, Orbit-Stabilizer]
    We have the following divisibility results:
    \begin{itemize} 
        \item Let $H$ be a subgroup of $G$, let $[G:H]$ denote the number of cosets of $H$ in $G$, then 
        \begin{equation*}
            |G|=|H|[G:H]
        \end{equation*}
        \item Let $G$ be a finite group acting transitively on a finite set $A$, then for any $a\in A$, we have 
        \begin{equation*}
            |\text{Stab}_G(a)|\cdot|O_G(a)|=|G|
        \end{equation*}
    \end{itemize} 
\end{prop}
The class formula is when $G$ acts on itself by conjugation:
\begin{prop}[class formula]
    Let $G$ act on a finite set $S$, and let $Z$ denote fixed points of this action, then 
    \begin{equation*}
        |S|=|Z|+\sum_{a\in A}|O_G(a)|
    \end{equation*}
    where $A$ includes exactly one element from each nontrivial orbit.

    If $G$ acts on itself by conjugation, then 
    \begin{equation*}
        |G|=|Z(G)|+\sum_{g}|[g]|=|Z(G)|+\sum_{g}\frac{|G|}{|C_G(g)|}
    \end{equation*}
    where $[g]$ denote the conjugacy class of $g$, and the sum includes exactly one from each nontrivial conjugacy class in $G$.
\end{prop}


\begin{prob}[F2019-Q2]
    2. Let \( p, q \) be two prime numbers such that \( p \mid q - 1 \). Prove that  
    \begin{itemize}
    \item[(a)]there exists an integer \( r \neq 1 \mod q \) such that \( r^p \equiv 1 \mod q \);
    \item[(b)] there exists (up to an isomorphism) only one noncommutative group of order \( pq \).
    \end{itemize}
\end{prob}
\begin{proof}
    \begin{itemize}
        \item[(a)] We want to show that there exists an element $r\in(\Z/q\Z)^\times$ such that 
        \begin{equation*}
            r^p\equiv 1\mod q
        \end{equation*}
        We can do this because $(\Z/q\Z)^\times$ has order $(q-1)$ and $p\vert (q-1)$. Therefore by Cauchy's theorem, there exists an element of order $p$ in $(\Z/p\Z)^\times$.
        \item[(b)] Let $n_p,n_q$ denote the number of $p$, $q$-Sylow subgroups. We see that $n_q\vert p$ and $n_q\equiv 1\mod q$, since $p<q$, we must have $n_q=1$. Now $n_p=1$ or $q$ by the same reasoning. Suppose $n_q=1$, let $P,Q$ denote the normal subgroups of order $p,q$, then
        \begin{equation*}
            G\cong P\times Q
        \end{equation*}
        by a standard argument (included in the lemma below). Then $G$ is commutative. Since $G$ is noncommutative, we have $n_p=q$. Choose any $p$-Sylow subgroup $P$, we know that 
        \begin{equation*}
            G\cong Q\rtimes_\theta P
        \end{equation*} 
        where $Q$ is the normal subgroup of order $q$ and $\theta: P\to\text{Aut}(Q)=(\Z/q\Z)^\times$. We know either $\theta: 1\mapsto 1$, is the trivial map which produces  a commutative group; or $\theta: 1\mapsto r$, where $r\in(\Z/q\Z)^\times$ is some element of order $p$. 
        
        
        % By (a), such an element exists. Since $(q-1)$ is even, we know there only exists $1$ element in $(\Z/q\Z)^\times=\Z/(q-1)\Z$ of order $2$, namely, $r$. Therefore 
        % \begin{equation*}
        %     G\cong Q\rtimes_\theta P
        % \end{equation*}
        % where $\theta: P\to\text{Aut}(Q)=(\Z/q\Z)^\times$ is such that $\theta: 1\mapsto r$ defines the noncommutative group of order $pq$ up to isomorphisms.
    \end{itemize}
\end{proof}



\begin{lem}
    Let $p,q$ be two primes such that $q\nmid (p-1)$, and $N$, $H$ has order $p,q$ respectively, suppose that $N$ is normal in $G$, and $N\cap H=\{e\}$, then 
    \begin{equation*}
        G\cong N\times H
    \end{equation*}
\end{lem}
\begin{proof}
    We consider the map 
    \begin{equation*}
        \psi: N\times H\to G
    \end{equation*}
    such that 
    \begin{equation*}
        (n,h)\mapsto nh
    \end{equation*}
    We want to show that $\psi$ is a homomorphism and $\psi$ is injective (hence bijective by size argument). It is clearly injective: 
    \begin{equation*}
        nh=e\Rightarrow n, h\in N\cap H=\{e\}
    \end{equation*}
    It suffices to show that $\psi$ is a homomorphism. We see that this implies 
    \begin{equation*}
        n_1n_2h_1h_2=n_1h_1n_2h_2
    \end{equation*}
    Therefore it suffices to for any $n\in N, h\in H$, one has
    \begin{equation*}
        nh=hn
    \end{equation*}
    Consider the conjugation action 
    \begin{equation*}
        \varphi: H\to \text{Aut}(N)
    \end{equation*}
    where 
    \begin{equation*}
        h\mapsto \left( n\mapsto hnh^{-1}\right)
    \end{equation*}
    Then we claim that $\varphi$ is trivial. This is because $\ker(\varphi)$ has size either $1$ or $q$. If it has size $q$, then the map is trivial; if it has size $1$, then $H$ embeds in $\text{Aut}(N)$, however, $|H|=q, \text{Aut}(N)=p-1$, and $q\nmid(p-1)$, hence impossible. This shows that the map is trivial, i.e., for $n\in N, h\in H$, 
    \begin{equation*}
        hn=nh
    \end{equation*}
    as desired. 
\end{proof}


\begin{prob}[F2015-Q1]
    Prove every group of order $15$ is cyclic.
\end{prob}
\begin{proof}
    We will show that any group $G$ of order $15$ is isomorphic to 
    \begin{equation*}
        G\cong\frac{\Z}{3\Z}\times\frac{\Z}{5\Z}
    \end{equation*}
    For this, using the above lemma, it suffices to show that there is one normal subgroup of order $3$ and one normal subgroup of order $5$. We repeat the argument above, $n_5\mid 3$ and $n_5\equiv 1\mod 5$, hence $n_5=1$. Moreover, $n_3\mid 5$ and $n_3\equiv 1\mod 3$, hence $n_3=1$ as well. By the lemma above, we know that 
    \begin{equation*}
        G\cong\frac{\Z}{3\Z}\times\frac{\Z}{5\Z}
    \end{equation*}
    hence cyclic as desired.
\end{proof}

\begin{prob}[S2013-Q2]
Let \( p \) and \( q \) be primes with \( p < q \). Let \( G \) be a group of order \( pq \). Prove the following statements:

\begin{itemize}
    \item[(a)] If \( p \) does not divide \( q - 1 \) (i.e., \( p \nmid q - 1 \)), then \( G \) is cyclic.

    \item[(b)] If \( p \) divides \( q - 1 \) (i.e., \( p \mid q - 1 \)), then \( G \) is either cyclic or isomorphic to a non-abelian group on two generators. Give the presentation of this non-abelian group.
\end{itemize}
\end{prob}
\begin{proof}
    This question is exactly the same as  F19-Q2, we will only outline here.
    \begin{itemize}
        \item[(a)] We have $n_q=1$, and $n_p\mid q$, hence $n_p=1$ or $q$, moreover $n_p\equiv 1\mod p$. If $n_p=q$, this implies that $p\mid(q-1)$, hence $n_p=1$. Therefore by the above argument
        \begin{equation*}
            G\cong\frac{\Z}{p\Z}\times\frac{\Z}{q\Z}
        \end{equation*}
        \item[(b)] If $p\mid(q-1)$, then $n_p=1$ or $q$. Hence $G$ is either of the form above or isomorphic to the non-abelian group 
        \begin{equation*}
            G=Q\rtimes_\theta P
        \end{equation*}
        We know from F2019-Q2, the trivial $\theta$ defines the abelian, hence cyclic group $G=\frac{\Z}{p\Z}\times\frac{\Z}{q\Z}$. And $\theta: 1\mapsto r$, for some $r\in(\Z/q\Z)^\times$ of order $p$ defines a non-abelian group. 

    \textcolor{red}{not finished, what are the two generators}        
    \end{itemize}
\end{proof}

\begin{prob}[F2007-Q1]
    Prove that no group of order $148$ is simple.
\end{prob}
\begin{proof}
    We note the prime factorization of $148$ is 
    \begin{equation*}
        148=2^2\cdot 37
    \end{equation*}
    We see that $n_{37}\mid 4$ and $n_{37}\equiv 1\mod 37$, therefore $n_{37}=1$. This shows that there exists a normal subgroup of order $37$, i.e., the group is not simple.
\end{proof}


\begin{prob}[F2017-Q1]
    Show that there is no simple group of order $30$.
\end{prob}
\begin{proof}
    This is slightly more complicated, and we will use a counting argument. 
    Same reasoning as the above. The prime factorization of $30$ is as below:
    \begin{equation*}
        30=2\cdot 3\cdot 5
    \end{equation*}
    We see $n_5\mid 6$, and $n_5\equiv 1\mod 5$. Unfortunately, $n_5$ could either be $1$ or $6$. Now $n_3\mid 10$, and $n_3\equiv 1\mod 3$, unfortunately again $n_3$ could be $10$. However, we argue that $n_3=10$ and $n_5=6$ cannot happen at the same time. Suppose this is the case, then there are $20$ elements of order $2$ and $24$ elements of order $5$, but this is too many! Hence either $n_3=1$ or $n_5=1$, as desired.
\end{proof}




\begin{prob}[F2011-Q1]
    \phantom{}
        \begin{itemize}
            \item[(a)] Let \( G \) be a group of order 5046. Show that \( G \) cannot be a simple group. You may not appeal to the classification of finite simple groups.
            
            \item[(b)] Let \( p \) and \( q \) be prime numbers. Show that any group of order \( p^2q \) is solvable.
        \end{itemize}   
\end{prob}
\begin{proof} 
    The proof is very similar like above.
    \begin{itemize}
        \item[(a)] The prime factorization of $5049$ is as follows:
        \begin{equation*}
            5049=2\cdot 3\cdot 29^2
        \end{equation*}
        Hence we see $n_{29}=1$, i.e., there is a normal subgroup of order $29$, therefore not simple.
        \item[(b)] We will do discussion by cases.
        \begin{itemize}
            \item[(1)] $p>q$. Then $n_p=1$ or $q$ and $n_p\equiv 1\mod p$, therefore $n_p=1$. Let $P$ be the normal subgroup of $G$ of order $p^2$, we thus have 
            \begin{equation*}
                \{e\}\subset P\subset G
            \end{equation*}
            It is clear that $|G/P|=q$, thus abelian, and $|P|=p^2$ also abelian as well (by the lemma below). This shows that $G$ is solvable.
            \item[(2)] $p<q$. Then $n_p=1$ or $q$, and $n_q=1$ or $p^2$. Suppose that $n_q=1$, let $Q$ denote the normal subgroup of order $q$, then 
            \begin{equation*}
                \{e\}\subset Q\subset G
            \end{equation*}
            It is clear that $Q$ and $G/Q$ are both abelian. Suppose that $n_q=p^2$ instead, then there are only $p^2q-p^2(q-1)=p^2$ elements of order $\neq q$. Since any $p$-Sylow subgroup has $p^2$ elements with order $\neq q$, we must have $n_p=1$. 
            Hence we are in case (1) again. This shows that $G$ is solvable in either case $n_q=1, p^2$.
        \end{itemize}
    \end{itemize}
\end{proof}


\begin{lem}[$p^2$ abelian]
    Fix prime $p$, any group of order $p^2$ is abelian. 
\end{lem}
\begin{proof}
    For any nontrivial $p$ group, by the class formula, the center $Z(G)$ is nontrivial, thus the center has order either $p$ or $p^2$. If it has order $p^2$, then the group is abelian. If it has order $p$, then 
    \begin{equation*}
        \left|G/Z(G)\right|=p
    \end{equation*}
    is also cyclic, therefore $G$ is abelian (strictly speaking is a contradiction that $|Z(G)|=p$). In either case, we see that $G$ is abelian.
\end{proof}

\begin{prob}
    Any $p$-group is solvable, for any prime $p$. 
\end{prob}
\begin{proof}
    Suppose $|G|=p^r$ for some $r\geq 0$, we will use induction on $r$. If $r=0$, then the trivial group is trivally solvable. 
    \begin{itemize}
        \item Base case: if $r=1$, $|G|=p$, then $G$ is cyclic, hence solvable. 
        \item Induction step: suppose that $G$ is solvable for all $|G|=p^k$, where $0\leq k\leq r-1$. Now we want to show that $G$ of order $p^r$ is solvable. We know $G$ has a nontrivial center, suppose that $|Z(G)|=p^k$, where $1\leq k\leq r$, then 
        \begin{equation*}
            |G/Z(G)|=p^{r-k}, 0\leq r-k\leq r-1 
        \end{equation*}
        We know any group $G$ is solvable if and only if there exists a sequence of subgroups $G_0,\dots,G_k$
        \begin{equation*}
            \{e\}=G_0\subset\dots\subset G_k=G
        \end{equation*}
        such that $G_{i-1}$ is normal in $G_i$ and $G_i/G_{i-1}$ is solvable. Therefore we see when $|G|=p^{r}$, 
        \begin{equation*}
            \{e\}\subset Z(G)\subset G
        \end{equation*}
        has $Z(G)$ solvable, and $G/Z(G)$ also solvable by the induction hypothesis, so we close the induction. 
    \end{itemize}
    

    % We know that any $p$-group $G$ has nontrivial center $Z(G)$. Suppose that $|G|=p^r$ for some $r$. If $|Z(G)|=p^r$, then $G$ is solvable. If $|Z(G)|=p^{r-1}$, then 
    %  Then we also have $G$ is solvable. Suppose that $|Z(G)|=p^{r-2}$, then $|Z(G)|$ is 
\end{proof}

\begin{prob}[S2016-Q1]
    Classify all groups of order $66$, up to isomorphism.
\end{prob}
\begin{proof}
    By $66=2\cdot 3\cdot 11$, we know $n_{11}=1$. We claim that there is a normal subgroup isomorphic to $\Z/33\Z$. 
    \begin{enumerate}
        \item First we show that there is a subgroup of order $33$. Let $P_{11}$ denote the normal subgroup of order $11$ and let $P_{3}$ denote a $3$-Sylow subgroup of $G$. Then we claim that the following
        \begin{equation*}
            H=\{gh: g\in P_{11}, h\in P_{3}\}
        \end{equation*}
        forms a subgroup and is isomorphic to $\Z/33\Z$. By the Lemma 1.1, we see that 
        \begin{equation*}
            H\cong\frac{\Z}{11\Z}\times\frac{\Z}{3\Z}=\frac{\Z}{33\Z}
        \end{equation*}
        \item Now we show that it is normal. This follows from the following general lemma:
        \begin{lem}
            Let $p$ be the smallest prime factor of $|G|$, and let $H$ be a subgroup with index $p$, then $H$ is normal.
        \end{lem}
        \begin{proof}
            We will only prove in the case that $H$ is a subgroup of index $2$, i.e., $G=H\sqcup (G\setminus H)$. We see for all $g\in G$, 
            \begin{equation*}
                gH=Hg
            \end{equation*}
            since if $g\in H$, then the equality holds; if $g\not\in H$, then $gH=G\setminus H$, so is $Hg$.  
        \end{proof}
        Now since there is a subgroup of order $2$, we can write $G$ as a semidirect product 
        \begin{equation*}
            G=\frac{\Z}{33\Z}\rtimes_\theta \frac{\Z}{2\Z}
        \end{equation*}
        The number of nonisomorphic groups will depend on the choice of $\theta$. There are four different choices for $\theta: H\to\text{Aut}\left(\frac{\Z}{11\Z}\times\frac{\Z}{3\Z}\right)=\frac{\Z}{10\Z}\times\frac{\Z}{2\Z}$
        \begin{equation*}
            \begin{cases}
                \theta_1: 1\mapsto (0,0)\\
                \theta_2: 1\mapsto (0,1)\\
                \theta_3: 1\mapsto (5, 0)\\
                \theta_4: 1\mapsto (5,1)
            \end{cases}
        \end{equation*}
    \end{enumerate} 
    \textcolor{red}{what is happening}
\end{proof}



\begin{prob}[S2007-Q2]
    Prove that no group of order $224$ is simple.
\end{prob}
\begin{proof}
    The prime factorization is 
    \begin{equation*}
        224=2^5\cdot 7
    \end{equation*}
    If $n_2=1$ or $n_7=1$, then we are done; assume that $n_2=7$ instead, then we recall $G$ has a nontrivial transitive action on the set of $2$-Sylow subgroups, i.e., there is a homomorphism $\varphi:G\to S_7$. We know $\ker(\varphi)$ is a normal subgroup of $G$. Since the action is nontrivial transitive, we know $\ker(\varphi)\neq G$. If $\ker(\varphi)=\{e\}$, then $\varphi$ produces an embedding of $G$ into $S_7$. However, $|G|=224\nmid |S_7|$. This shows that $\ker(\varphi)$ is a nontrivial proper normal subgroup of $G$, concluding that $G$ is not simple.
\end{proof}

\begin{prob}[F2008-Q1]
    Show that no group of order $36$ is simple.
\end{prob}
\begin{proof}
    \begin{equation*}
        36=2^2\cdot 3^3
    \end{equation*}
    We know $n_2\mid 9, n_2\equiv 1\mod 2$, and $n_3\mid 4, n_3\equiv 1\mod 3$. We know $n_3=1$ or $4$, suppose that $n_3=4$, then there is a nontrivial action of $G$ on the set of $3$-Sylow subgroups, i.e., 
    \begin{equation*}
        \varphi: G\to S_4
    \end{equation*}
    Suppose that $G$ is simple, we know $\ker(\varphi)\neq G$ since the action is nontrivial, by assumption $\ker(\varphi)=\{e\}$, which implies that $\varphi$ is an embedding, but $|G|=32\nmid |S_4|$, which is a contradiction. This implies that $G$ is not simple.
\end{proof}

\begin{prob}[S2014-Q2]
    All groups of order less than $60$ are solvable, i.e., there exists a sequence of subgroups of $G$, $G_0, \dots, G_k$ such that $G_i$ is normal in $G_{i+1}$ and $G_{i+1}/G_i$ is abelian, and 
    \begin{equation*}
        1=G_0\subset\dots\subset G_k=G
    \end{equation*}
\end{prob}
\begin{proof}
    Groups of order $p, pq, p^2, p^2q$ are solvable.
    \begin{align*}
        \big\{ 1, 2, 3, 4, 5, 6, 7, 9, 10, 11, 12, 13, 14, 15, 17, 19, 20, 21, 22, 23, 25, 26, 28, 29, 30, \\ 31, 33, 34, 35, 37, 38, 39, 41, 43, 44, 45, 46, 47, 49, 50, 51, 52, 53, 55, 57, 58, 59 \big\}
    \end{align*}
    And any $p$-group is also solvable. 
    \begin{equation*}
        \{8, 16, 27, 32\}
    \end{equation*}
    The remaining groups are 
    \begin{equation*}
        \left\{24, 36, 40, 42, 48, 54, 56 \right\}
    \end{equation*}
    \begin{itemize}
        \item[24:] If $n_2=1$ or $n_3=1$, then we are done. We see $n_2=1$ or $3$, consider the action $\varphi:G\to S_3$. We see $\ker(\varphi)$ is a proper normal subgroup of $G$, this implies that 
        \begin{equation*}
            \{e\}\subset \ker(\varphi)\subset G
        \end{equation*}
        where $|\ker(\varphi)|$ is a known solvable group, hence we are done.
        \item[36:] Exactly same as above, we assume $n_3\neq 1$, therefore $n_3=4$, the action $\varphi: G\to S_4$ is not injective, hence $\ker(\varphi)$ is again a proper normal subgroup of $G$ that is solvable.
        \item[40:] We see $n_5=1$, therefore 
        \begin{equation*}
            \{e\}\subset \Z/5\Z\subset G
        \end{equation*}
        \item[42:] We see $n_7=1$.
        \item[48:] We see $n_2=1$ or $3$, the the action $\varphi: G\to S_3$ is not injective, hence $\ker(\varphi)$ is a proper normal subgroup of $G$ that is solvable.
        \item[54:] We see $n_3=1$. 
        \item[56:] We know $n_7=1$ or $8$ and $n_2=1$ or $7$. The group action argument does not work. We assume $n_7=8$, then there can be at most $56-8(7-1)=8$ elements of order $\neq 7$. This shows that $n_2=1$. Hence 
        \begin{equation*}
            \{e\}\subset P_2\subset G
        \end{equation*}
    \end{itemize}
\end{proof}

\begin{prob}[S2012-Q1]
    Let \( G \) be a group of order \( p^3 q^2 \), where \( p \) and \( q \) are prime integers. Show that for \( p \) sufficiently large and \( q \) fixed, \( G \) contains a normal subgroup other than \(\{1\}\) and \( G \).
\end{prob}
\begin{proof}
    We want to show that there exists a normal group of size $p^3$, i.e., $n_p=1$. We know $n_p\mid q^2, n_p\equiv 1\mod p$. Let $p$ be large enough such that $p>(q^2-1)$, then ths forces $n_p=1$, as desired.
\end{proof}


% Semidirect product: $N$ is normal, then there is an action of $H$ on $\text{Aut}(N)$, if $|H|=p$, then this map has kernel either be trivial or the entire group, in the case that it is trivial, this map is an embedding.

% Same argument for the simple group question.


\begin{prob}[F2014-Q4]
    \phantom{text}
    \begin{itemize}
        \item[(a)] Let \( G \) be a group of order \( p^2q^2 \), where \( p \) and \( q \) are distinct odd primes, with \( p > q \). Show that \( G \) has a normal subgroup of order \( p^2 \).
        
        \item[(b)] Can a solvable group contain a non-solvable subgroup? Explain.
    \end{itemize}
\end{prob}
\begin{proof}
    \begin{itemize}
        \item[(a)] We know $n_p=1$ or $q$ or $q^2$, and $n_p\equiv 1\mod p$. Since $p>q$, we know $n_p\neq q$. It suffices to show that $n_p\neq q^2$: suppose that $n_p=q^2$, then 
        \begin{equation*}
            p\mid (q^2-1)=(q+1)(q-1)
        \end{equation*}
        Since $p$ is prime, $p\mid (q+1)$ or $p\mid(q-1)$. The latter impossible since $q<p$. $p\mid (q+1)$ is also impossible because this implies that $q=p+1$, which implies that $q$ is even, a contradiction.
        \item[(b)] It is not possible. Suppose $G$ is a solvable group, let $H$ be a subgroup of $G$, then we know there exists sequence 
        \begin{equation*}
            \{e\}=G_0\subset\dots\subset G_k=G
        \end{equation*}
        such that $G_i$ is normal in $G_{i+1}$ and $\frac{G_{i+1}}{G_i}$ is abelian. We define $H_i=G_i\cap H$, then we see $H$ is solvable with sequence $H_0\subset\dots H_k$.
    \end{itemize}

\end{proof}


\begin{prob}[F2018-Q2]
    Let \( G \) be a group of order \( 24 \). Assume that no Sylow subgroup of \( G \) is normal in \( G \). Show that \( G \) is isomorphic to the symmetric group \( S_4 \).
\end{prob}
\begin{proof}
    
\end{proof}


\begin{prob}[F2001-Q1]
    Let \( G \) be a finite group and let \( N \) be a normal subgroup of \( G \) such that \( N \) and \( G/N \) have relatively prime orders.
        
        \begin{enumerate}
            \item Assume that there exists a subgroup \( H \) of \( G \) having the same order as \( G/N \). Show that \( G = HN \). (Here \( HN \) denotes the set \(\{xy \mid x \in H, y \in N\}\).)
            
            \item Show that \( \phi(N) = N \), for all automorphisms \( \phi \) of \( G \).
        \end{enumerate}
\end{prob}
\begin{proof}
    \begin{enumerate}
        \item Since $N,H$ have relatively prime orders, $N\cap H=\{e\}$, thus we can write 
        \begin{equation*}
            G=N\rtimes_\theta H
        \end{equation*}
        where $\theta(h)n=hnh^{-1}$. One can show that the map $\varphi: N\rtimes_\theta H\to G$ as
        \begin{equation*}
            \varphi:(n,h)\mapsto nh
        \end{equation*}
        It is clear that $\varphi$ is a homomorphism and injective, thus by a size argument we have $\varphi$ is an isomorphism. This shows $G=NH$ and similarly $G=HN$.
        \item Any automorphism $\phi$ of $G$ permutes the $p$-Sylow subgroups. Suppose that $|G|=p_1^{i_1}\dots p_k^{i_k}$, then after rearranging,
        \begin{equation*}
            |N|=p_1^{i_1}\dots p_j^{i_j}
        \end{equation*}
        because $N$ and $G/N$ have relatively prime orders. Hence $N$ contains all the Sylow $p_i$-subgroups, hence $\phi(N)=N$ for all automorphisms $\phi$ of $G$.
    \end{enumerate}
\end{proof}


\begin{prob}[S2001-Q1]
    Let \( G \) be a finite group and \( p \) the smallest prime number dividing the order \( |G| \) of \( G \). Let \( H \) be a subgroup of \( G \) of index \( p \) in \( G \). Show that \( H \) is necessarily a normal subgroup of \( G \).
\end{prob}
\begin{proof}
    $G$ has an action on $G/H$ by left multiplication: $\varphi: G\to\text{Aut}(G/H)$ such that 
    \begin{equation*}
        \varphi(g)(\bar{g}H)=g\bar{g}H
    \end{equation*}
    We will show that $H=\ker(\varphi)$. First we see that $\ker(\varphi)\subset H$: 
    \begin{equation*}
        \ker(\varphi)=\{g\in G: g\bar{g}H=\bar{g}H: \text{ for all }\bar{g}\in G\}
    \end{equation*}
    letting $\bar{g}\in H$ we see $g\in\ker(\varphi)$ implies $g\in H$, i.e., $\ker(\varphi)\subset H$. 

    Now we use a size argument to show $|H|\leq|ker\varphi|$. We note that $\im(\varphi)$ is a subgroup of $\text{Aut}(G/H)=S_p$, thus 
    \begin{equation*}
        \frac{|G|}{|\ker(\varphi)|}\text{ divides } p!
    \end{equation*}
    because $\frac{|G|}{|\ker(\varphi)|}$ also divides $|G|$ and $p$ is the smallest prime that divides $p$, we must have 
    \begin{equation*}
        \frac{|G|}{|\ker(\varphi)|} \text{ divides } p
    \end{equation*}
    Note that $\frac{|G|}{|H|}=p$, this gives 
    \begin{equation*}
        |H|\leq|\ker(\varphi)|
    \end{equation*}
    which shows $H\subset\ker(\varphi)$, hence $H=\ker(\varphi)$.
\end{proof}




\begin{center}\textcolor{pink}{(End of Page 5)}
\end{center}

\section{Class Formula, Classification of $p$-groups}
\begin{defn}[nilpotent group]
    Let $G$ be a group. Define inductively an increasing sequence 
$\{e\}=Z_0 \subseteq Z_1 \subseteq Z_2 \subseteq \cdots$ of subgroups of $G$ 
as follows: for $i \geq 1$, $Z_i$ is the subgroup of $G$ corresponding 
to the center of $G/Z_{i-1}$. One can show that $Z_i$ is normal in $G$. A group is \emph{nilpotent} if $Z_m = G$ for some $m$.
\end{defn}
\begin{example}
    \phantom{text}
    \begin{itemize}
        \item $p$-groups are nilpotent.
        \item Nilpotent groups are solvable.
    \end{itemize}
\end{example}
\begin{prop}
    We have the following classification of groups of order $p, p^2, p^3$, for prime $p$.
    \begin{itemize}
        \item $|G|=p$ implies $G\cong\Z/p\Z$.
        \item $|G|=p^2$ implies 
        \begin{equation*}
            G\cong\frac{\Z}{p^2\Z} \quad\text{ or }\quad G\cong\frac{\Z}{p\Z}\oplus\frac{\Z}{p\Z}
        \end{equation*}
        \item $|G|=p^3$ implies that 
        \begin{equation*}
            G\cong\frac{\Z}{p^3\Z} \quad \text{ or } \quad G/Z(G)\cong\frac{\Z}{p\Z}\oplus\frac{\Z}{p\Z} \quad\text{ or }\quad [G,G]=Z(G)
        \end{equation*}
    \end{itemize}
\end{prop}
% \begin{proof}
%     \textcolor{red}{finish proof}
% \end{proof}

\begin{prob}[S2010-Q1]
    Let \( G \) be a non-abelian group of order \( p^3 \), where \( p \) is prime. Determine the number of distinct conjugacy classes in \( G \).
\end{prob}
\begin{proof}
    We know $G$ has a nontrivial center, and if $|Z(G)|=p^2$ or $p^3$, then $G$ is abelian, this shows that $|Z(G)|=p$, now let $g\in G\setminus Z(G)$, then 
    \begin{equation*}
        Z(G)\subsetneq Z_g(G)\subsetneq G
    \end{equation*}
    where $Z(G)\subsetneq Z_g(G)$ because $g\in Z_g(G)$, and $Z_g(G)\subsetneq G$ since $g\not\in Z(G)$. This shows that $Z_g(G)$ is a subgroup of order $p^2$, in other words, the size of the conjugacy class of any $g\in G\setminus Z(G)$ is 
    \begin{equation*}
        |[g]|=\left|\frac{G}{Z_g(G)}\right|=p
    \end{equation*}
    By the class formula, 
    \begin{equation*}
        |G|=|Z(G)|+\sum_{a\in A}|[a]|
    \end{equation*}
    where $A$ contains one $a$ from each nontrivial conjugacy class $[a]$. Thus we have 
    \begin{equation*}
        p^3=p+Np\Rightarrow N=p^2-1
    \end{equation*}
    Every element in $Z(G)$ is its own conjugacy class, thus the total number of conjugacy classes is 
    \begin{equation*}
        p^2+p-1
    \end{equation*}
\end{proof}

\begin{prob}[F2013-Q1]
    Let $p>2$ be a prime. Classify groups of order $p^3$ up to isomorphism. The two nonabelian groups of order \( p^3 \) (for \( p \neq 2 \)), up to isomorphism, are:

    \[
    \text{Heis}(\mathbb{Z}/(p)) = 
    \left\{ 
    \left( \begin{array}{ccc}
    1 & a & b \\
    0 & 1 & c \\
    0 & 0 & 1
    \end{array} \right) \middle| a, b, c \in \mathbb{Z}/(p)
    \right\}
    \]
    
    and
    
    \[
    G_p = 
    \left\{ 
    \left( \begin{array}{cc}
    a & b \\
    0 & 1
    \end{array} \right) \middle| a, b \in \mathbb{Z}/(p^2), a \equiv 1 \bmod p
    \right\}
    \]
    
    \[
    = 
    \left\{ 
    \left( \begin{array}{cc}
    1 + pm & b \\
    0 & 1
    \end{array} \right) \middle| m, b \in \mathbb{Z}/(p^2)
    \right\}
    \]
\end{prob}

\begin{prob}[F2014-Q5]
    \phantom{text}
    \begin{enumerate}
        \item[(a)] Prove that every group of order \( p^2 \) (with \( p \) prime) is abelian. Then classify such groups up to isomorphism.
    
        \item[(b)] Give an example of a non-abelian group of order \( p^3 \) for \( p = 3 \). \\
        \emph{Suggestion: Represent the group as a group of matrices.}
    \end{enumerate}    
\end{prob}
\begin{proof}
    \begin{itemize}
        \item[(a)] See Lemma 1.2. There are two abeliean groups: $\frac{\Z}{p^2\Z}, \frac{\Z}{p\Z}\times\frac{\Z}{p\Z}$.
        \item[(b)] See Problem 1.18.
    \end{itemize}
\end{proof}

\begin{prob}[F2019-Q4, S2015-Q3]
    Find all irreducible representations of a finite $p$-group over a field of characteristic $p$.
\end{prob}
\begin{proof}
    Let $G$ any finite $p$-group. Let $V$ be an irreducible representation over $\F_p$, consider the $[\F_pG]$-module $W$ generated by any $v\in V\setminus\{0\}$. We see $W$ is a finite-dimensional vector space over $\F_p$, i.e., 
    \begin{equation*}
        |W|=p^d
    \end{equation*}
    for some $d\geq 1$. We consider the action of $G$ on $W$, all the orbits of this action either has size $1$ or is a power of $p$, since $G$ is a $p$-group, by the class formula, let $N$ be the number of nontrivial orbits of size $1$, 
    \begin{equation*}
        |W|\equiv 1+N\mod p\Rightarrow 1+N\equiv 0\mod p
    \end{equation*}
    Hence there exists at least one nontrivial orbit $\{v\}$ of size $1$. We consider the vector space $\bar{W}$ generated by $v$ over $\F_p$: it is one-dimensional vector space contained in $V$, invariant under $G$, since $V$ is irreducible, we must have $V=\bar{W}$. The action of $G$ on $\bar{W}$ is the trivial action, thus all irreducible representations of a finite $p$-group over $\F_p$ are trivial.
\end{proof}


\section{Random Problems}

\begin{prob}[F2010-Q1]
    Let \( G \) be a group. Let \( H \) be a subset of \( G \) that is closed under group multiplication. Assume that \( g^2 \in H \) for all \( g \in G \). Show that:
    \begin{itemize}
        \item \( H \) is a normal subgroup of \( G \)
        \item \( G/H \) is abelian
    \end{itemize}
\end{prob}
\begin{proof}
    \begin{itemize}
        \item We first show that $H$ is subgroup. It remains to show that if $h\in H$, then $h^{-1}\in H$, we know $(h^{-1})^2\in H$, thus 
        \begin{equation*}
            h(h^{-1})^2=h^{-1}\in H
        \end{equation*}
        as desired.  Now we show that $H$ is normal: for any $h\in H$, $g\in G$, we want to show $ghg^{-1}\in H$.
        \begin{align*}
            ghg^{-1}&=(gh)^2(gh)^{-1}hg^{-1}\\
            &=(gh)^2h^{-1}g^{-1}hg^{-1}\\
            &=(gh)^2h^{-1}(g^{-1}h)^2(g^{-1}h)^{-1}g^{-1}\\
            &=(gh)^2h^{-1}(g^{-1}h)^2h^{-1}\in H
        \end{align*}
        as desired.
        \item It suffices to show that for any $g_1,g_2\in G$, we have 
        \begin{equation*}
            g_1g_2H\subset g_2g_1H
        \end{equation*}
        Take any $h\in H$, we want to show $(g_2g_1)^{-1}g_1g_2h\in H$,
        \begin{align*}
            (g_2g_1)^{-1}g_1g_2h&=(g_2g_1)^{-2}g_2g_1^2g_2h\\
            &=(g_2g_1)^{-2}(g_2g_1^2)^2(g_2g_1^2)^{-1}g_2h\\
            &=(g_2g_1)^{-2}(g_2g_1^2)^2g_1^{-2}h\in H
        \end{align*}
        as desired.
    \end{itemize}
\end{proof}

\begin{prob}[S2014-Q1]
    Find the number of colorings of the faces of a cube using 3 colors, where two colorings are considered equal if they can be transformed into each other by a rotation of the cube.
    
    [\textit{Hint}: Use Burnside's formula:
    
    \[
    |X/G| = \frac{1}{|G|} \sum_{g \in G} |X^g|,
    \]
    
    where a group $G$ acts on a set $X$, $X/G$ is the set of orbits, and for every $g \in G$, $X^g$ is the fixed subset of $g$ in $X$.]
\end{prob}
\begin{proof}
    Let $X$ be the set of all possible colorings of the cube (equal cubes allowed), we have $|X|=3^6$. We notice two things:
    \begin{enumerate}
        \item The group of rotations of a cube is $S_4$.
        \item For $\sigma_1,\sigma_2\in S_4$ that are conjugates of each other, $|X^{\sigma_1}|=|X^{\sigma_2}$. Therefore for the Burnside's formula becomes 
        \begin{equation*}
            |X/S_4|=\frac{1}{|S_4|}\sum_{[\sigma]
            \text{ conj classes}}|[\sigma]|\cdot|X^\sigma|
        \end{equation*}
    \end{enumerate}
    Now we analyze for each conjugacy class $[\sigma]$, what is $|X^\sigma|$.
    \begin{itemize}
        \item $(1+1+1+1)$, $|[e]|=1$ and $|X^e|=3^6$.
        \item $(1+1+2)$, $|[\sigma_1]|=6$ and $|X^{\sigma_1}|=3^3$.
        \item $(1+3)$, $|[\sigma_2]|=8$, and $|X^{\sigma_2}|=3^2$.
        \item $(2+2)$, $|[\sigma_3]|=6$, and $|X^{\sigma_3}|=3^4$.
        \item $(4)$, $|[\sigma_4]|=6$, and and $|X^{\sigma_4}|=3^3$.
    \end{itemize}
    Thus combining we get 
    \begin{equation*}
        |X/S_4|=\frac{1}{24}\left(3^6+6\cdot 3^3+8\cdot 3^2+6\cdot 3^4+6\cdot 3^3\right)=57
    \end{equation*}
\end{proof}

\begin{prob}[S2019-Q4]
    Let \( f \) be a polynomial with \( n \) variables and define
\[
\operatorname{Sym}(f) = \{ \sigma \in S_n \mid f(x_{\sigma(1)}, x_{\sigma(2)}, \ldots, x_{\sigma(n)}) = f(x_1, x_2, \ldots, x_n) \}.
\]

\begin{enumerate}
    \item Prove that \(\operatorname{Sym}(f)\) is a subgroup of \( S_n \).
    
    \item Prove that the dihedral group \( D_4 \) (the group of symmetries of the square) is isomorphic to \(\operatorname{Sym}(x_1 x_2 + x_3 x_4)\).
\end{enumerate}
\end{prob}
\begin{proof}
    \begin{enumerate}
        \item The group $S_n$ acts on the polynomial ring $k[x_1,\dots,x_n]$, by permuting the $x_i$ to $x_{\sigma(i)}$, and we see that Sym$(f)$ is the centralizer of a fixed element $f\in k[x_1,\dots,x_n]$, hence is a subgroup.
        \item We have a total of $8$ elements in Sym$(x_1x_2+x_3x_4)$: 
        \begin{equation*}
            \{e, (12), (34), (12)(34), (13)(24), (14)(23), (1324), (1423)\}
        \end{equation*}
        and we can by drawing a square tha this corresponds to the group $D_4$.
    \end{enumerate}
\end{proof}

\begin{prob}[S2011-Q1, F2004-Q1]
    \phantom{text}
    \begin{itemize}
    \item[(a)] Let \( H \) be a proper nontrivial subgroup of a finite group \( G \) (i.e., \( H \neq \{1\} \) and \( H \neq G \)).  
    Prove that \( G \) is not the union of all conjugates of \( H \) in \( G \).
    
    \item[(b)] Give an example of an infinite group \( G \) for which the assertion in part (a) fails.
    \end{itemize}
\end{prob}
\begin{proof}
    \begin{itemize}
        \item[(a)] If $H$ is normal, then all conjugations of $H$ is equal to $H$, but $H\subsetneq G$, this $G$ is not not the union of all conjugates of \( H \) in \( G \). Now suppose $H$ is not normal, assume the contrary that $G$ is the union of all conjugates of $H$, then the number of distinct conjugates of $H$ is $[G:N_G(H)]$, hence 
        \begin{equation*}
            |G|=[G:N_G(H)]\cdot|H|\iff [G:H]=[G:N_G(H)]\iff [N_G(H):H]=1
        \end{equation*}
        this is a contradiction since $H$ is not normal. Thus $G$ not the union of all conjugates of \( H \) in \( G \).
        \item[(b)] Consider 
        \begin{equation*}
            B=\left\{\begin{pmatrix}
                *&*\\
                0&*
            \end{pmatrix}\right\}\subset \text{GL}_2(\mathbb{C})
        \end{equation*}
        It is clear that conjugation of matrices in $B$ do not give matrices with nonzero left bottom entry.
    \end{itemize}
\end{proof}

\begin{prob}[S2009-Q1]
    Let \( H \) and \( K \) be two solvable subgroups of a group \( G \) such that \( G = HK \).
    
    \begin{enumerate}
        \item Show that if either \( H \) or \( K \) is normal in \( G \), then \( G \) is solvable.
        
        \item Give an example where \( G \) may not be solvable without the assumption in (a).
    \end{enumerate}
\end{prob}
\begin{proof}
    \begin{enumerate}
        \item WLOG suppose that $H$ is normal, then the composite map $\varphi=\pi\circ\iota$: 
        \[\begin{tikzcd}
            K & G & {G/H}
            \arrow["\iota", from=1-1, to=1-2]
            \arrow["\pi", from=1-2, to=1-3]
        \end{tikzcd}\]
        is surjective, therefore 
        \begin{equation*}
            \{e\}\subset H\subset G
        \end{equation*}
        $G/H\cong K/\ker(\varphi)$ is solvable, hence $G$ is solvable.
        \item The smallest nonsolvable group is $A_5$, we have 
        \begin{equation*}
            A_5=HK
        \end{equation*}
        where $H=\la(12345)\ra, K=A_4=\{\sigma\in A_5: \sigma(5)=5\}$. Now $H, K$ are both solvable, but $G$ is not.
    \end{enumerate}
\end{proof}

\begin{prob}[F2003-Q1]
    In a group \( G \), let \( 1 \) denote the identity element and let \([x, y] = xyx^{-1}y^{-1}\) denote the commutator of elements \(x, y \in G\).
    
    \begin{enumerate}
        \item Express \([z, xy]x\) in terms of \(x\), \([z, x]\), and \([z, y]\).
        
        \item Prove that if the identity \([[x, y], z] = 1\) holds in \(G\), then the following identities hold in \(G\):
        \[
        [x, yz] = [x, y][x, z] \quad \text{and} \quad [xy, z] = [x, z][y, z].
        \]
    \end{enumerate}
\end{prob}
\begin{proof}
    \begin{enumerate}
        \item We have 
        \begin{align*}
            [z,xy]x&=zxyz^{-1}y^{-1}x^{-1}x\\
            &=zxz^{-1}x^{-1}xzyz^{-1}y^{-1}\\
            &=[z,x]x[z,y]
        \end{align*}
        \item The identity $[[x,y],z]=1$ implies 
        \begin{equation*}
            [x,y]z=z[x,y]
        \end{equation*}
        Therefore using the identity in 1, we have 
        \begin{align*}
            [x,yz]&=[x,y]y[x,z]y^{-1}\\
            &=[x,y]yy^{-1}[x,z]\\
            &=[x,y][x,z]
        \end{align*}
        Similarly
        \begin{align*}
            [xy,z]&=xyzy^{-1}x^{-1}z^{-1}\\
            &=xyzy^{-1}z^{-1}zx^{-1}z^{-1}\\
            &=x[y,z]x^{-1}[x,z]\\
            &=[y,z][x,z]\\
            &=[x,z][y,z]
        \end{align*}
    \end{enumerate}
\end{proof}

\begin{prob}[S2005-Q1]
    Let $k$ be a field. Let $G = \mathrm{GL}_n(k)$ be the general linear group, where $n > 0$. Let $D$ be the subgroup of diagonal matrices, and let $N = N_G(D)$ be the normalizer of $D$ in $G$. Determine the quotient group $N/D$.
\end{prob}

\begin{prob}[F2009-Q1]
    Let \( G \) be a finite group, and let \( \mathrm{Aut}(G) \) be its automorphism group. Consider the group action \( \phi \colon \mathrm{Aut}(G) \times G \to G \) defined by \( \phi(\sigma, g) = \sigma(g) \). Assume \( G \) has exactly two orbits under this action.
    
    \begin{enumerate}
        \item Determine all such groups \( G \) up to isomorphism.
        
        \item For each case from (a), determine when \( \mathrm{Aut}(G) \) is solvable.
    \end{enumerate}
\end{prob}

\begin{prob}[F2016-Q1]
    Determine $\text{Aut}(S_3)$.
\end{prob}
\begin{proof}
    Every element $\sigma\in\text{Aut}(S_3)$ must send order $2$ elements $\{(12), (23), (13)\}$ to one another, and order $3$ elements $\{(123), (132)\}$ to each other. However, $\sigma$ is determined by how it permutes 
    \begin{equation*}
        \{(12), (23), (13)\}
    \end{equation*}
    Thus every $\sigma$ is an inner automorphism of the form $\sigma_g(h)=ghg^{-1}$ for $g,h\in S_3$ and $g$ is some transposition. Hence 
    \begin{equation*}
        \text{Aut}(S_3)\cong S_3
    \end{equation*}
\end{proof}


\chapter{Representation Theory}


% \begin{thm}[Maschke's theorem]

% \end{thm}

% \begin{lem}[Schur's Lemma]
    
% \end{lem}

\begin{prop}[properties of characters]
    
\end{prop}
\begin{prop}
    The character tables for $S_3, S_4, A_5, S_5$ are as follows:
\end{prop}

\begin{thm}[Compilation of theorems]
    Schur's lemma: 
    \begin{enumerate}
        \item If $\varphi: V\to W$ is a $G$-invariant map, i.e., 
        \begin{equation*}
            \varphi(\rho(g))=\rho(g)\varphi(v)
        \end{equation*} 
        where $V,W$ are irreducible representations, then $\varphi=0$ or an isomorphism. This is true for any field $k$ that $V,W$ are over.
        \item If $\varphi: V\to V$ and everything as above, then 
        \begin{equation*}
            \varphi(v)=\lambda v
        \end{equation*}
        for some $\lambda\in k^\times$. This is only true when $k$ is algebraically clsoed.
        \item $\text{Hom}_G(V,W)\begin{cases}
            k \text{ if } V\cong W\\
            0 \text{ if  not}
        \end{cases}$, where $V,W$ are irreducible. This is true for $k$ algebraically closed.
        \item Mascheke's theorem: any finite dimensional representation $V$ of a finite group $G$ can be decomposed into a direct sum of irreducible representations.
        \begin{equation*}
            V=V_1^{r_1}\oplus\dots\oplus V_k^{r_k}
        \end{equation*}
        where $V_i$'s are irreducible. This is true when the characteristic $k$ does not divide $|G|$, notably this always holds for characteristic $0$ fields.
        \item Do not mix them up.
    \end{enumerate}
\end{thm}

\section{Characters}

\begin{prob}[S2008-Q4]
    Let \( V \cong \mathbb{C}^n \) be an \( n \)-dimensional complex vector space with standard basis \( e_1, \ldots, e_n \). Consider the permutation action \( S_n \times V \to V \) defined by:
    \[
    \sigma \cdot e_i = e_{\sigma(i)} \quad \text{for } \sigma \in S_n
    \]
    Decompose \( V \) into irreducible \( \mathbb{C}[S_n] \)-modules.
\end{prob}

\begin{prob}[S2014-Q5]
    Find the table of characters for $S_4$.
\end{prob}

\begin{prob}[F2016-Q6]
    Find a table of characters for the alternating group $A_5$.
\end{prob}


\begin{prob}[F2015-Q3]
    Let \( G = S_4 \) (the symmetric group on four letters).

\begin{enumerate}
    \item[(a)] Prove that \( G \) has two non-equivalent irreducible complex representations of dimension 3; call them \( \rho_1 \) and \( \rho_2 \).
    
    \item[(b)] Decompose the tensor product representation \( \rho_1 \otimes \rho_2 \) into a direct sum of irreducible representations of \( G \).
\end{enumerate}
\end{prob}


\begin{prob}[F2011-Q4]
    Let \(\rho \colon S_3 \to \mathrm{GL}(2, \mathbb{C})\) be a two-dimensional irreducible representation of the symmetric group \(S_3\). 

\begin{enumerate}
    \item Decompose the tensor square \(\rho^{\otimes 2}\) into irreducible representations of \(S_3\).
    \item Decompose the tensor cube \(\rho^{\otimes 3}\) into irreducible representations of \(S_3\).
\end{enumerate}
\end{prob}

\begin{prob}[F2014-Q3]
    Let \( G = S_3 \) be the symmetric group on three elements.

\begin{enumerate}
    \item[(a)] Prove that \( G \) has an irreducible complex representation of dimension 2 (call it \( \rho \)), but none of higher dimension.
    
    \item[(b)] Decompose the triple tensor product \( \rho \otimes \rho \otimes \rho \) into a direct sum of irreducible representations of \( G \).
\end{enumerate}
\end{prob}

\begin{prob}[S2006-Q6]
    Let \( S_4 \) be the symmetric group on four elements.
    
    \begin{enumerate}
        \item Give an example of a non-trivial 8-dimensional complex representation of \( S_4 \).
        
        \item Show that every 8-dimensional complex representation of \( S_4 \) contains a 2-dimensional invariant subspace.
    \end{enumerate}
\end{prob}

\begin{prob}[F2007-Q5]
    Prove that every 5-dimensional complex representation of the alternating group \( A_4 \) (the alternating group of degree 4) contains a 1-dimensional invariant subspace.
\end{prob}

\begin{prob}[S2004-Q6]
    Consider complex representations of a finite group \( G \). Let \( \sigma_1, \ldots, \sigma_s \) be representatives of the conjugacy classes of \( G \), and let \( \chi_1, \ldots, \chi_s \) be the irreducible characters of \( G \).
    
    \begin{enumerate}
        \item Define an inner product on the \(\mathbb{C}\)-vector space of class functions on \( G \) such that \( \{\chi_1, \ldots, \chi_s\} \) forms an orthonormal basis.
        
        \item Let \( A = (a_{ij}) \) be the character table matrix of \( G \), where \( a_{ij} = \chi_i(\sigma_j) \) for \( 1 \leq i, j \leq s \). Prove that \( A \) is invertible.
    \end{enumerate}
\end{prob}

\begin{prob}[S2018-Q4, S2007-Q5]
    Is $S_4$ isomorphic to a subgroup of $\text{GL}_2(\mathbb{C})$?
\end{prob}

\begin{prob}[S2010-Q6]
    Let \( G \) be a group of order 24. Using representation theory, prove that \( G \neq [G, G] \), where \([G, G]\) denotes the commutator subgroup of \( G \).
\end{prob}
\begin{proof}
    Suppose $G=[G,G]$, then we claim the only $1$-dimensional representation $\rho:G\to\mathbb{C}^\times$ is the trivial one. This is because if $\rho$ is one-dim, then
    \begin{equation*}
        [G,G]\subset\ker(\rho)
    \end{equation*}
    i.e., $\rho$ is trivial. However, there is no way to write 
    \begin{equation*}
        |G|=24=1+d_1^2+\dots+d_k^2
    \end{equation*}
    where $d_i\geq 2$. Thus $G\neq [G,G]$.
\end{proof}

\begin{prob}[F2017-Q6]
    Let \( G \) be a finite group with center \( Z(G) \). Show that if \( G \) admits a faithful irreducible representation \( \rho \colon G \to \mathrm{GL}_n(k) \) for some positive integer \( n \in \mathbb{Z}^+ \) and some field \( k \), then the center \( Z(G) \) is cyclic.
\end{prob}
\begin{proof}
    (We will only do the case where $k$ is algebraically closed). For any $z\in Z(G)$, $\rho(z): V\to V$ is a $G$-map, i.e., 
    \begin{equation*}
        \rho(z)(\rho(g)v)=\rho(g)(\rho(z)v)
    \end{equation*}
    We know by Schur's lemma that $\rho(z)$ is a scalar multiplication: 
    \begin{equation*}
        \rho(z)\in k^\times
    \end{equation*}
    Because $\rho$ is faithful, $Z$ embeds into $k^\times$ via $\rho$. 
    \begin{lem}[Fact]
        Any finite subgroup of $k^\times$ for field $k$ is cyclic.
    \end{lem}
    Hence $Z$ is cyclic.
\end{proof}


\begin{prob}[S2005-Q6]
    Let \( V \) be a finite-dimensional vector space over a field \( k \), and let \( G \) be a finite group with an irreducible representation \( \varphi \colon G \to \mathrm{GL}(V) \). Suppose \( H \) is a finite abelian subgroup of \( \mathrm{GL}(V) \) contained in the centralizer of \( \varphi(G) \). Prove that \( H \) must be cyclic.
\end{prob}



\begin{prob}[F2010-Q6]
    Let \( G \) be a non-abelian group of order \( p^3 \), where \( p \) is prime.
    
    \begin{enumerate}
        \item Determine the number of isomorphism classes of irreducible complex representations of \( G \), and find their dimensions.
        
        \item Which of these irreducible complex representations are faithful? Justify your answer.
    \end{enumerate}
\end{prob}


\begin{prob}[S2011-Q5]
    Let \( K \) be a field, and let \( \Phi \colon G \to \mathrm{GL}_n(K) \) be an \( n \)-dimensional matrix representation of a group \( G \). Define a \( G \)-action on the matrix ring \( M_n(K) \) by:
\[
(g, A) \mapsto \Phi(g) \cdot A \quad \text{(matrix multiplication)}
\]
for \( g \in G \) and \( A \in M_n(K) \). This action induces a group homomorphism \( \Psi \colon G \to \mathrm{GL}(M_n(K)) \). Express the character \( \chi_\Psi \) of \( \Psi \) in terms of \( \chi_\Phi \) (the character of \( \Phi \)).
\end{prob}

\begin{prob}[S2015-Q5]
    Prove that a tensor product of irreducible representations over an algebraically closed field is irreducible.
\end{prob}

\begin{prob}[S2001-Q3]
    Calculate the complete character table for $\Z/3\Z\times S_3$, where $S_3$ is the symmetric group in 3 letters.
\end{prob}

\section{Induced representations}


\begin{prob}[S2009-Q6]
    Let \( G = S_4 \) and consider the subgroup \( H = \langle (1\,2), (3\,4) \rangle \).
    
    \begin{enumerate}
        \item[(a)] Determine the number of irreducible complex characters of \( H \).
        
        \item[(b)] Choose a non-trivial irreducible character \( \psi \) of \( H \) over \( \mathbb{C} \) satisfying \( \psi((1\,2)(3\,4)) = -1 \). Compute the values of the induced character \( \operatorname{ind}_H^G(\psi) \) on all conjugacy classes of \( G \), and express it as a sum of irreducible characters of \( G \).
    \end{enumerate}
\end{prob}

\section{Frobenius Reciprocity}

\begin{prob}[S2017-Q6]
    Let $G$ be a finite group and H an abelian subgroup. Show that every
    irreducible representation of $G$ over $\mathbb{C}$ has dimension 
    $\leq[G : H]$.
\end{prob}

\begin{prob}[S2008-Q6]
    Give an example of non-isomorphic finite groups with same character
table. Construct the character table in detail.
\end{prob}

\begin{prob}[S2012-Q4]
    Let \( Q \) be the quaternion group with presentation:
\[
Q = \langle t, s_i, s_j, s_k \mid t^2 = 1,\ s_i^2 = s_j^2 = s_k^2 = s_i s_j s_k = t \rangle.
\]

\begin{enumerate}
    \item[(a)] Find four non-isomorphic 1-dimensional real representations of \( Q \).
    
    \item[(b)] Prove that the natural embedding \( \rho \colon Q \to \mathbb{H} \) given by:
    \[
    \rho(t) = -1, \quad \rho(s_i) = i, \quad \rho(s_j) = j, \quad \rho(s_k) = k
    \]
    defines an irreducible 4-dimensional real representation of \( Q \), where \( \mathbb{H} \) is the algebra of real quaternions.
    
    \item[(c)] Classify all irreducible complex representations of \( Q \) up to isomorphism.
\end{enumerate}
\end{prob}

\begin{prob}[F2004-Q6]
    Let \( D_8 \) be the dihedral group of order 8, with presentation:
    \[
    D_8 = \langle r, s \mid r^4 = 1 = s^2,\ rs = sr^{-1} \rangle.
    \]
    
    \begin{enumerate}
        \item Determine all conjugacy classes of \( D_8 \).
        
        \item Find the commutator subgroup \( D_8' \) of \( D_8 \) and determine the number of distinct degree-1 (linear) characters of \( D_8 \).
        
        \item Construct the complete complex character table of \( D_8 \).
    \end{enumerate}
\end{prob}

\begin{prob}[F2000-Q7]
    Let \( D_{10} \) be the dihedral group of order 10, with presentation:
    \[
    D_{10} = \langle r, s \mid r^5 = 1 = s^2,\ rs = sr^{-1} \rangle.
    \]
    
    \begin{enumerate}
        \item Determine all conjugacy classes of \( D_{10} \).
        
        \item Compute the commutator subgroup \( D_{10}' \) of \( D_{10} \).
        
        \item Prove that \( D_{10}/D_{10}' \cong \mathbb{Z}/2\mathbb{Z} \) and deduce that \( D_{10} \) has exactly two distinct degree-1 characters.
        
        \item Construct the complete complex character table of \( D_{10} \).
    \end{enumerate}
\end{prob}

























\chapter{Semisimple Algebra}

% Page 19-20

% The most recent semisimple question appeared in 2020.

% (2020, 2017, 2016, 2011, 2010, 2008, 2006)

\begin{defn}[Division ring]
    Any nonzero element in a unit.
\end{defn}
\begin{prop}
    Let $A$ be a semisimple finite-dimensional algebra over $F$, then $A$ can be decomposed into a direct sum of matrix algebras over a division ring:
    \begin{equation*}
        A=M_{n_1}(D_1)\oplus\dots\oplus M_{n_k}(D_k)
    \end{equation*}
    where $D_i$'s are division rings, $M_{n_i}$ is the algebra of $n_i\times n_i$ matrices with entries in $D_i$. This decomposition is unique up to permutation.

    For example, let $G$ be a finite, group, then the group algebra $\C(G)$ can be decomposed to 
    \begin{equation*}
        \C(G)=M_{n_1}(\C)\oplus\dots\oplus M_{n_k}(\C)
    \end{equation*}
    where $|G|=\sum_id_i^2$, and $k$ is the number of conjugacy classes of $G$. Hence it suffices to compute the irreducible representations of $G$.
\end{prop}

\begin{prop}
    Any semisimple ring \( R \) can be decomposed into a finite direct sum of simple ideals \( J_i \):
\[
R = \bigoplus_{i=1}^n J_i.
\]
\end{prop}

\begin{prob}[F2019-Q5]
    Determine the number of two-sided ideals in the group algebra \(\mathbb{C}[S_3]\), where \(S_3\) is the symmetric group of permutations of \(\{1, 2, 3\}\).
\end{prob}
\begin{proof}
    Using the Proposition above, we know that 
    \begin{equation*}
        \C(S_3)=M_1(\C)+M_1(\C)+M_2(\C)=\C\oplus\C\oplus M_2(\C)
    \end{equation*}
\end{proof}



\begin{prob}[F2009-Q6, F2001-Q5]
    Let \( \rho \colon G \to \mathrm{GL}_n(\mathbb{C}) \) be an irreducible complex representation of a finite group \( G \), with character \( \chi \), and let \( C \) be the center of \( G \).
    
    \begin{enumerate}
        \item Prove that for every \( s \in C \), the matrix \( \rho(s) \) is a scalar multiple of the identity matrix \( I_n \).
        
        \item Using part (a), show that \( |\chi(s)| = n \) for all \( s \in C \).
        
        \item Establish the inequality \( n^2 \leq [G : C] \), where \( [G : C] \) is the index of \( C \) in \( G \).
        
        \item Prove that if \( \rho \) is faithful (i.e., injective), then \( C \) must be cyclic.
    \end{enumerate}
\end{prob}



\begin{proof}
    \begin{enumerate}
        \item $\mathbb{C}$ is algebraically closed therefore Schur's lemma applies (see F2017-Q6)
        \item We know that 
        \begin{equation*}
            \rho(z)=\begin{pmatrix}
                \lambda&0&\dots&0\\
                0&\lambda&\dots&0\\
                \dots\\
                0&\dots&0&\lambda
            \end{pmatrix}
        \end{equation*}
        We also know that $C$ is finite and $\rho(z^r)=I$, which implies $|r|=1$. This gives $|\chi(s)|=n$ for all $s\in C$. 
        \item We know that $\rho$ is irreducible, hence the corresponding character $\chi$ satisfies 
        \begin{equation*}
            \frac{1}{|G|}\sum_{g\in G}|\chi(g)|^2=\frac{|C|}{|G|}n^2+\frac{1}{|G|}\sum_{g\not\in C}|\chi(g)|^2=1
        \end{equation*}
        This implies that 
        \begin{equation*}
            \frac{|C|}{|G|}n^2\leq 1\Rightarrow n^2\leq[G:C]
        \end{equation*}
        \item If $\rho$ is faithful, then $C$ embeds into $k^\times$, and any finite subgroup of $k^\times$ is cyclic.
    \end{enumerate}
\end{proof}

\begin{prob}[S2017-Q5]
    Prove directly from the definition of (left) semisimple ring that every such
ring is (left) Noetherian and Artinian. (You may freely use facts about
semisimple, Noetherian, and Artinian modules.)
\end{prob}
\begin{proof}
    Any semisimple ring $R$ can be decomposed into a finite direct sum of simple ideals $I_i$:
    \begin{equation*}
        R=\bigoplus_{i=1}^nI_i
    \end{equation*}
    This directly implies that the ascending and descending chain condition: there aren't infinitely sequence of ideals of strict inclusions.
\end{proof}





\begin{prob}[S2005-Q4]
    Let \( R \) be a ring and \( L \) a minimal left ideal of \( R \) (i.e., \( L \) contains no non-zero proper left ideals of \( R \)). Assuming \( L^2 \neq 0 \), prove that \( L = Re \) for some non-zero idempotent element \( e \in R \).
\end{prob}
\begin{proof}
    We recall that a ring element $e\in R$ is an idempotent if and only if 
    \begin{equation*}
        e^2=e
    \end{equation*}
    It suffices to show that there exists a nonzero idempotent element $e\in L$ since $Re$ is an ideal contained in $L$, thus $Re=L$. Take any $x\neq 0$ in $L$, such that there exists $g\in L$ such that $gx\neq 0$ (this is guaranteed by $L^2\neq 0$). The ideal $Lx$ is contained in $L$, since $L$ is simple, we must have 
    \begin{equation*}
        L=Lx
    \end{equation*}
    Hence $x\in L$ can be written as 
    \begin{equation*}
        x=ex
    \end{equation*}
    for some $e\in L$, multiplying both sides by $e$ and moving terms, we get 
    \begin{equation*}
        (e^2-e)x=0
    \end{equation*}
    It suffices to show that 
    \begin{equation*}
        \left\{g\in L: gx=0\right\}=\{0\}
    \end{equation*}
    This is because $\{g\in L: gx=0\}$ is again an ideal contained in $L$, since we assumed that there exists some $g\in L$ such that $gx\neq 0$, 
    \begin{equation*}
        \{g\in L:gx=0\}=\{0\}
    \end{equation*}
    and we are done!
\end{proof}


\begin{prob}[S2016-Q6, F2006-Q6, F2008-Q6]
    Let \( A \) be a finite-dimensional semisimple algebra over \( \mathbb{C} \), and let \( V \) be an \( A \)-module that decomposes as \( V \cong S \oplus S \), where \( S \) is a simple \( A \)-module. Determine the automorphism group \( \operatorname{Aut}_A(V) \) of \( V \) as an \( A \)-module.
\end{prob}
\begin{proof}
   By Schur's lemma, since $S$ is a simple $A$-module, we know 
   \begin{equation*}
        \text{End}_A(S)\cong\C
   \end{equation*}
   Thus 
   \begin{equation*}
        \text{End}(V)\cong M_2(\C)
   \end{equation*}
   hence 
   \begin{equation*}
        \text{Aut}_A(V)\cong \text{GL}_2(\C)
   \end{equation*}
\end{proof}



\begin{prob}[S2010-Q5]
    Classify all non-commutative semi-simple rings with $512$ elements.
(You can use the fact that finite division rings are fields.)
\end{prob}
\begin{proof}
    By Artin-Wedderburn, we know that this finite semisimple ring can be decomposed into a finite direct sum of matrix rings:
    \begin{equation*}
        R\cong M_{n_1}(F_1)\oplus\dots\oplus M_{n_k}(F_k)
    \end{equation*}
    where $F_i$ are finite fields. Further more we can assume $n_1\geq n_2\geq\dots\geq n_k$.The total number of elements is 
    \begin{equation*}
        F_1^{n_1^2}\dots F_k^{n_k^2}=512=2^9
    \end{equation*}
    Thus we see all the components are powers of $2$. Since $R$ is noncommutative, we may assume that $n_1\geq 2$. If $n_1=3$, then $F_1=\F_2$, we have 
    \begin{equation*}
        R\cong M_{3}(\F_2)
    \end{equation*}
    If $n_1=2$, we can have $n_2=2$, then 
    \begin{equation*}
        R\cong M_2(\F_2)\oplus M_2(\F_2)\oplus \F_2
    \end{equation*}
    or $n_2=n_3=\dots=n_k=1$, then we have (different ways of adding to $5$):
    \begin{equation*}
        R\cong M_2(\F_2)\oplus\begin{cases}
            \F_2\oplus\F_2\oplus\F_2\oplus\F_2\oplus F_2\\
            \F_{4}\oplus \F_2\oplus\F_2\oplus F_2\\
            \F_8\oplus\F_2\oplus\F_2\\
            \F_{16}\oplus\F_2\\
            \F_8\oplus\F_4\\
            \F_{16}\oplus F_2\\
            \F_{32}
        \end{cases}
    \end{equation*}



    % There are three semi-simple rings of of order $512$ that are not commutative. They are 
    % \begin{equation*}
    %     \F_2\oplus\F_2\oplus\F_2\oplus\F_2\oplus\F_2\oplus M(\F_4)
    % \end{equation*},
    % \begin{equation*}
    %     \F_2\oplus M(\F_4)\oplus M(\F_4)
    % \end{equation*}
    % and 
    
    % \begin{equation*}
    %     M(\F_9)
    % \end{equation*}
\end{proof}

\begin{prob}[F2011-Q5]
    Let \( A \) be a finite-dimensional semisimple algebra over \( \mathbb{C} \), and let \( V \) be a finitely-generated \( A \)-module. Prove that \( V \) has only finitely many \( A \)-submodules if and only if \( V \) decomposes into a direct sum of pairwise irreducible non-isomorphic (i.e., simple)  \( A \)-modules.
\end{prob}
\begin{proof}
    Suppose that $V$ is a direct sum of distinct irreducible $A$-modules, then 
    \begin{equation*}
        V=S_1\oplus\dots\oplus S_n
    \end{equation*}
    where $S_i$'s are nonisomorphic and simple. Hence the only submodules of $S_i$ is $\{0\}$ and $S_i$, i.e., there are only finitely many submodules of $V$.

    Conversely, we suppose that there are finitely many $A$-submodules of $V$, because $V$ is semisimple, we know 
    \begin{equation*}
        V=\bigoplus_{i=1}^n S_i^{n_i}
    \end{equation*} 
    where $S_i$'s are semisimple. It suffices to show that $n_i=1$ for all $i$. Suppose that 
    \begin{equation*}
        V=S_i\oplus S_i
    \end{equation*}
    By Schur's lemma, we have 
    \begin{equation*}
        \text{Hom}(S_i,S_i)\cong\C
    \end{equation*}
    there are infinitely many distinct $\phi:S_i\to S_i$, and we note that 
    \begin{equation*}
        \left\{(s,\phi(s)):\phi\in\text{Hom}(S_i,S_i)\right\}
    \end{equation*} 
    is a submodule of $V$, thus there are infinitely many submodules, which is a contradiction.
\end{proof}






















\chapter{Linear Algebra I}

Topics: finitely generated modules/PID, triangularization, diagonalization, Jordan canonical form.

% Page 21-23

\begin{prop}
    Assume that characteristic a linear operator $T:V\to V$ factors completely over $k$, then $T$ is diagonalizable if and only if the minimal polynomial splits into distinct linear factors (has no repeated roots).
\end{prop}


\begin{prob}[F2018-Q1]
    Let \( V \) be an \( n \)-dimensional vector space over a field \( k \) and let \( \alpha: V \to V \) be a linear endomorphism.  
Prove that the minimal and characteristic polynomials of \( \alpha \) coincide if and only if there is a vector \( v \in V \) so that:  

\[\{v, \alpha(v), \ldots, \alpha^{n-1}(v)\}\]
is a basis for \( V \).
\end{prob}
\begin{proof}
    % We will show there is no vector $v\in V$ such that $\{v, \alpha(v), \ldots, \alpha^{n-1}(v)\}$ is linearly independent if and only if the minimal and characteristic polynomials aren't the same. 
    % \begin{align*}
    %     \text{no $v$ s.t. $\{v,\dots, \alpha^{n-1}(v)\}$ is linearly independent}&\iff \text{ for all } v\in V, \{v,\ldots, \alpha^{n-1}v\} \text{is linearly dependent}\\
    %     &\iff \alpha^{n-1}v=c_0v+\dots+c_{n-2}\alpha^{n-2}v \text{ for all } v
    % \end{align*}
    % We then claim that $\alpha^{n-1}v=c_0v+\dots+c_{n-2}\alpha^{n-2}v \text{ for all } v$ if and only if the minimal polynomial $p_m$ has degree less than the characteristic polynomial $p_n$. Note that $p_n$ has degree $n$, and we can write 
    % \begin{equation*}
    %     p_n(t)=a_nt^n+\dots+a_1t+a_0
    % \end{equation*}
    % and $p_n(\alpha)v=0$ for all $v\in V$, therefore $\alpha^{n-1}v=c_0v+\dots+c_{n-2}\alpha^{n-2}v \text{ for all } v$ if and only if 
    % \begin{equation*}
    %     p_n(\alpha)v=a_n\alpha^{n-1}v+\dots+a_1\alpha v+a_0=p_{n-1}(\alpha)v=0
    % \end{equation*}
    % where $p_{n-1}$ is a polynomial of degree at most $n-1$ and is such that $p_{n-1}(\alpha)=0$. Thus no $v$ such that $\{v,\dots,\alpha^{n-v}v\}$ is linearly independent if and only if minimal and characteristic polynomial have different degrees.
\end{proof}


\begin{prob}[F2018-Q3]
    \phantom{text}
    \begin{itemize}
        \item[(a)] Fix a positive integer \( n \) and classify all finite modules over the ring \( \mathbb{Z}/n\Z \).
        \item[(b)] Prove, either using (a) or from first principles, for a fixed prime \( p \) that all finite modules over \( \mathbb{Z}/p\Z \) are free.
    \end{itemize}
\end{prob}
\begin{proof}
    By the classification of finite abelian groups
    \begin{itemize}
        \item[(a)] $G\cong\bigoplus_{i,j}\frac{\Z}{p_i^{ij}\Z}$, but $G$ cannot be a $\Z/n\Z$-module unless $p_i^{ij}$ divides $n$ for all $i$:
        \begin{equation*}
            reasoning
        \end{equation*}
        Thus for a fixed $n$, let $n=p_1^{a_1}\dots p_k^{a_k}$ be its prime factorization, then 
        \begin{equation*}
            G\cong \bigoplus_{p\mid n}\bigoplus_{i}\frac{\Z}{p^i\Z}
        \end{equation*}
        where $i\leq a_i$ for each $i$.
        \item[(b)] By (a).
    \end{itemize}
\end{proof}


\begin{prob}[F2017-Q2]
    Let \(\Lambda\) be a free abelian group of finite rank \(n\), and let \(\Lambda' \subset \Lambda\) be a subgroup of the same rank. Let \(x_1, \ldots, x_n\) be a \(\mathbb{Z}\)-basis for \(\Lambda\), and let \(x_1', \ldots, x_n'\) be a \(\mathbb{Z}\)-basis for \(\Lambda'\). For each \(i\), write \(x_i' = \sum_{j=1}^n a_{ij}x_j\), and let \(A := (a_{ij}) \in \operatorname{Mat}_{n \times n}(\mathbb{Z})\). Show that the index \([\Lambda : \Lambda']\) equals \(|\det A|\).
\end{prob}
\begin{proof}
    Up to some basis change, we can write 
    \begin{equation*}
        \Gamma'=d_1\Z\oplus\dots\oplus d_k\Z
    \end{equation*}
    given $\Gamma=\Z\oplus\dots\oplus\Z$. Then Since we are taking the determinant, it is invariant under change of basis. One can compute the matrix using the standard basis for $\Gamma$ and $\Gamma'$, and it is clear that $[\Gamma:\Gamma']=\prod_{i=1}^kd_i=|\det(A)|$.
\end{proof}


\begin{prob}[S2001-Q5]
    \phantom{text}
    \begin{itemize}
        \item[(a)] Prove that an \( n \times n \) matrix \( A \) with entries in the field \( \mathbb{C} \) of complex numbers, satisfying \( A^3 = A \), can be diagonalized over \( \mathbb{C} \).
        
        \item[(b)] Does the statement in (a) remain true if one replaces \( \mathbb{C} \) by an arbitrary algebraically closed field \( F \)? Why or why not?
    \end{itemize}
\end{prob}
\begin{proof}
    \begin{itemize}
        \item[(a)] $A$ is diagonalizable if and only if the minimal polynomial splits into distinct linear factors. The characteristic polynomial is $p(t)=t(t+1)(t-1)$ and the minimal polynomial $p_m\mid p$ thus $A$ is diagonalizable.
        \item[(b)] This is not true. Take $k$ to be a field of characteristic $2$, then 
        \begin{equation*}
            p(t)=t(t^2-1)=t(t-1)^2
        \end{equation*}
        Thus the minimal polynomial could be $(t-1)^2$, i.e., $A$ is not necessarily diagonalizable.
    \end{itemize}
\end{proof}



\begin{prob}[F2001-Q3]
    Let \( A \) be an \( n \times n \) complex matrix with \( A^m = 0 \) for some integer \( m > 0 \).
    
    \begin{enumerate}
        \item Show that if \( \lambda \) is an eigenvalue of \( A \), then \( \lambda = 0 \).
        
        \item Determine the characteristic polynomial of \( A \).
        
        \item Prove that \( A^n = 0 \).
        
        \item Construct a \( 5 \times 5 \) matrix \( B \) satisfying \( B^3 = 0 \) but \( B^2 \neq 0 \).
        
        \item For any \( 5 \times 5 \) complex matrix \( M \) with \( M^3 = 0 \) and \( M^2 \neq 0 \), is \( M \) necessarily similar to your matrix \( B \) from part (d)? Justify your answer.
    \end{enumerate}
\end{prob}
\begin{enumerate}
    \item Suppose $\lambda$ is an eigenvalue, then there exists $v\neq 0$, such that 
    \begin{equation*}
        A^mv=\lambda^mv=0\Rightarrow \lambda=0
    \end{equation*}
    \item The characteristic polynomial is $p(t)=t^n$.
    \item Cayley-Hamilton theorem.
    \item Can have 
    \begin{equation*}
        B=\begin{pmatrix}
            0&1&0&0&0\\
            0&0&1&0&0\\
            0&0&0&0&0\\
            0&0&0&0&0\\
            0&0&0&0&0
        \end{pmatrix}
    \end{equation*}
    The important is that the top left $3\times 3$ matrix $A$ satisfies $A^3=0, A^3\neq 0$. This is constructed by building $B$ using the Jordan form.
    \item No, the lower $2\times 2$ matrix could be 
    \begin{equation*}
        \begin{pmatrix}
            0&0\\
            0&0
        \end{pmatrix}\quad \text{ or }\quad \begin{pmatrix}
            0&1\\
            0&0
        \end{pmatrix}
    \end{equation*}
\end{enumerate}


\begin{prob}[F2018-Q4]
    In this question all modules are left modules.

Let \( k \) be a field of characteristic different from 2 and let \( G = \{e, g\} \) be the multiplicative group with two elements. Consider the group ring \( A = k[G] \).

\begin{enumerate}
    \item[(a)] Show that the \( A \)-module \( A \) is a direct sum of two ideals of \( A \).
    \begin{itemize}
        \item List all proper ideals of \( A \).
        \item Is \( A \) a principal ideal domain?
    \end{itemize}
    
    \item[(b)] Show that every \( A \)-module decomposes into a direct sum of simple \( A \)-modules.
    
    \item[(c)] Assume now that the characteristic of \( k \) is 2. Give an example of an \( A \)-module that cannot be decomposed into a direct sum of two simple \( A \)-modules.
\end{enumerate}
\end{prob}
\begin{proof}
    \textcolor{red}{not finished}
\end{proof}


\begin{prob}[S2003-Q3]
    Prove that if a linear operator on a complex vector space is diagonal in some basis, then its restriction to any invariant subspace \( L \) is also diagonal in some basis of \( L \).
\end{prob}
\begin{proof}
    The linear operator $T$ is diagonalizable if and only if the minimal polynomial has no repeated factors, i.e., 
    \begin{equation*}
        f_m(x)=(x-\lambda_1)\dots(x-\lambda_k)
    \end{equation*}
    And $T\vert_L$ has minimal polynomial dividng $f_m$, hence it also has no repeated factors, thus $T\vert_L$ is also diagonalizable.
\end{proof}


% \begin{center}
%     \textcolor{pink}{End of Page 22}
% \end{center}


\begin{prob}[S2017-Q4]
    Let \(M\) be an invertible \(n\times n\) matrix with entries in an algebraically closed field \(k\) of characteristic not 2. Show that \(M\) has a square root, i.e. there exists \(N\in\text{Mat}_{n\times n}(k)\) such that \(N^{2}=M\).
\end{prob}
\begin{proof}
    It suffices to show that every Jordan block 
    \[
        J_n(\lambda) = 
        \begin{bmatrix}
        \lambda & 1       & 0       & \cdots & 0 \\
        0       & \lambda & 1       & \cdots & 0 \\
        0       & 0       & \lambda & \ddots & \vdots \\
        \vdots  & \vdots  & \ddots  & \ddots & 1 \\
        0       & 0       & \cdots  & 0      & \lambda
        \end{bmatrix}
        \]
        where $\lambda\neq 0$ is a square. We will proceed using induction. When $n=2$, the square root of 
        \begin{equation*}
            \begin{bmatrix}
                \lambda&1\\
                0&\lambda
            \end{bmatrix}=\begin{bmatrix}
                \lambda^\frac{1}{2}&\frac{1}{2}\lambda^{-\frac{1}{2}}\\
                0&\lambda^\frac{1}{2}
            \end{bmatrix}^2
        \end{equation*} 
        Now assume that $J_{k}$ is a square up to $k=n-1$, we want to show $J_n$ also has a square root. We claim $J_n$ has the following square 
        \begin{equation*}
            J_n=\begin{bmatrix}
                B^2&x\\
                0&\lambda

            \end{bmatrix}=\begin{bmatrix}
                B&x\\
                0&\lambda^{1/2}
            \end{bmatrix}^2
        \end{equation*}
        where $B$ is a $(n-1)\times (n-1)$ matrix and $x=(x_1,\dots,x_{n-1}), 0=(0,\dots,0)$. It suffices to find such an $x$ exists. Let $b_1,\dots,b_{n-1}$ denote the row vectors of $B$, we must satisfy 
        \begin{equation*}
            \begin{cases}
                b_1\cdot x+x_1\lambda^\frac{1}{2}=0\\
                \dots\\
                b_{n-2}\cdot x+x_{n-2}\lambda^\frac{1}{2}=0\\
                b_{n-1}\cdot x+x_{n-1}\lambda^\frac{1}{2}=1\\
            \end{cases}
        \end{equation*}
        Namely, we need to find $x$ that satisfies 
        \begin{equation*}
            (B+\lambda^\frac{1}{2}I)x=\begin{bmatrix}
                0\\
                \dots\\
                0\\
                1
            \end{bmatrix}
        \end{equation*}
        Since $(B+\lambda^{1/2}I)$ is invertible, there exists a unique solution, hence such $x$ exsits, $J_n$ has a square root! 
\end{proof}



\begin{prob}[S2008-Q1]
    Let \(k\) be a field. Consider the subgroup \(B \subset \text{GL}_2(k)\) where
    \[B = \left\{ \begin{pmatrix} a & b \\ 0 & d \end{pmatrix} \mid a, b, d \in k, ad \neq 0 \right\}.\]
    \begin{itemize}
        \item[(a)] Let \(Z\) be the center of \(\text{GL}_2(k)\). Show that
        \[\bigcap_{x \in \text{GL}_2(k)} x^{-1} Bx = Z.\]
        \item[(b)] Assume \(k\) is algebraically closed. Show that
        \[\bigcup_{x \in \text{GL}_2(k)} x^{-1} Bx = \text{GL}_2(k).\]
        \item[(c)] Assume \(k\) is a finite field. Is the equation in (b) still  true?
    \end{itemize}
\end{prob}
\begin{proof}
    \begin{itemize}
        \item[(a)] Let $y\in \bigcap_{x \in \text{GL}_2(k)}$, then for all $x\in\text{GL}_2(k)$, we have $xyx^{-1}\in B$. This shows that 
        \begin{align*}
            xyx^{-1}\in B\text{ for all }x&\iff xyx^{-1}\begin{bmatrix}
                0\\
                1
            \end{bmatrix}\in\left\la\begin{bmatrix}
                0\\
                1
            \end{bmatrix}\right\ra\\
            &\iff x^{-1}\begin{bmatrix}
                0\\
                1
            \end{bmatrix}
            \text{ is a subspace for } y \text{ for all }x\\
            &\iff \text{ the whole vector space is the eigenspace of }y\\
            &\iff \text{ $y$ is a scalar}\\
            &\iff y\in Z  
        \end{align*}
        \item[(b)] If $k$ is algebraically closed, then any matrix can be written as a triangular matrix up to some basis change.
        \item[(c)] It's false for finite fields. (b) is true when only when the characteristic polynomial can be factored completely over $k$. 
        Take $k=\F_2$, then we know $x^2+x+1$ is irreducible over $\F_2$, then we notice the matrix 
        \begin{equation*}
            \begin{pmatrix}
                0&1\\
                1&1
            \end{pmatrix}
        \end{equation*}
        which has characteristic polynomial exactly this. This is a counterexample.

        (One can take $g\in\overline{\mathbb{F}_p}\setminus\mathbb{F}_p$, then the characteristic polynomial for the map of multiplication by $g:\overline{\mathbb{F}_p}\to \overline{\mathbb{F}_p}$ where $\overline{\mathbb{F}_p}=\mathbb{F}_{p^2}$ is a vector space over $\mathbb{F}$ the minimial polynomial is $(t-g)^2$ which is irreducible over $\mathbb{F}_p$.)
    \end{itemize}
\end{proof}

\begin{prob}[S2009-Q4]
    Let \(E\) be a finite-dimensional vector space over an algebraically closed field \(k\). Let \(A, B\) be \(k\)-endomorphisms of \(E\). Assume \(AB = BA\). Show that \(A\) and \(B\) have a common eigenvector.
\end{prob}
\begin{proof}
    Since $k$ is algebraically closed, we know there exists at least one eigenvector of $A$, i.e., there exists $\lambda$ such that $Av=\lambda v$ for some $v\neq 0$. We denote this eigenspace by $E_\lambda$, and we note that $E_\lambda$ is invariant under $B$: let $v\in E_\lambda$
    \begin{equation*}
        A(Bv)=\lambda(Bv)
    \end{equation*}
    thus $Bv\in E_\lambda$ as well. Then it suffices to find an eigenvector of $B$ living inside $E_\lambda$, this is done by noting $B\vert_{E_\lambda}$ has an eigenvector in $E_\lambda$, as desired.
\end{proof}


\begin{prob}[F2005-Q6]
    Let \(E\) be a finite-dimensional vector space over a field \(k\). Assume \(S, T \in \text{End}_k(E)\). Assume \(ST = TS\) and both of them are diagonalizable. Show that there exists a basis of \(E\) consisting of eigenvectors for both \(S\) and \(T\).
\end{prob}
\begin{proof}
    It is the same proof as above except now we do this for all $E_{\lambda_1},\dots, E_{\lambda_k}$.
\end{proof}

\begin{prob}[S2015-Q2]
    Let \(A,B\) be two commuting operators on a finite dimensional space \(V\) over \(\mathbb{C}\) such that \(A^{n}=B^{m}\) is the identity operator on \(V\) for some positive integers \(n,m\). Prove that \(V\) is a direct sum of 1-dimensional invariant subspaces with respect to \(A\) and \(B\) simultaneously.
\end{prob}
\begin{proof}
    Because 
    \begin{equation*}
        A^n=B^m=I
    \end{equation*}
    We know that the minimal polynomial of $A,B$ both have no repeated roots, because $(t^n-1), (t^m-1)$ factor completely over $\C$. This shows that $A,B$ are commuting diagonalizable matrices, thus they can be simultaneously diagonalized.
\end{proof}






\chapter{Linear Algebra II}
Topics: exterior power, tensor algebras, traces, determinants

% Page 24-25

\begin{prob}[F2016-Q5]
    Let A be a linear transformation of a finite dimensional vector space over a field of characteristic \(\neq 2\).
    \begin{itemize}
        \item[(1)] Define the wedge product linear transformation \(\wedge^{2}A=A\wedge A\).
        \item[(2)] Prove that
        \[tr(\wedge^{2}A)=\frac{1}{2}(tr(A)^{2}-tr(A^{2})).\]
    \end{itemize}
\end{prob}
\begin{proof}
    \begin{itemize}
        \item[(a)] We recall the wedge product of vector space $V\wedge V$ is given by the basis 
        \begin{equation*}
            \{v_i\wedge v_j: i<j\}
        \end{equation*}
        satisfying 
        \begin{equation*}
            v_i\wedge v_j=-v_j\wedge v_i
        \end{equation*}
        where $\{v_1,\dots, v_n\}$ is a basis for $V$. And we define 
        \begin{equation*}
            A\wedge A(v_i\wedge v_j)=A(v_i)\wedge A(v_j)
        \end{equation*}
        \item[(b)] Consider the matrix representation of $A=(A_{ij})$, on the basis $\{v_i\wedge v_j: i<j\}$, 
        \begin{align*}
            A\wedge A(v_i\wedge v_j)&=\sum_{k,l=1}^nA_{ki}A_{lj}(v_k\wedge v_l)\\
            &=\sum_{k<l}A_{ki}A_{lj}(v_k\wedge v_l)+\sum_{l<k}A_{ki}A_{lj}(v_k\wedge v_l)\\
            &=\sum_{k<l}A_{ki}A_{lj}(v_k\wedge v_l)-\sum_{l<k}A_{ki}A_{lj}(v_l\wedge v_k)
        \end{align*}
        Thus the diagonal term with respect to $v_i\wedge v_j$ is 
        \begin{equation*}
            A_{ii}A_{jj}-A_{ji}A_{ij}
        \end{equation*}
        Thus 
        \begin{equation*}
            \text{Tr}(A\wedge A)=\sum_{i<j}A_{ii}A_{jj}-A_{ji}A_{ij}
        \end{equation*}
        Now 
        \begin{equation*}
            \text{Tr}(A)^2=\sum_{i=1}^nA_{ii}^2+2\sum_{i<j}A_{ii}A_{jj}
        \end{equation*}
        and 
        \begin{align*}
            \text{Tr}(A^2)&=\sum_{k,l=1}^nA_{lk}A_{kl}\\
            &=\sum_{i=1}^nA_{ii}^2+2\sum_{k<l}A_{lk}A_{kl}
        \end{align*}
        Thus we see that 
        \[tr(\wedge^{2}A)=\frac{1}{2}(tr(A)^{2}-tr(A^{2}))\]
    \end{itemize}
\end{proof}


\begin{prob}[S2006-Q5]
    Let \(V\) be a finite-dimensional vector space over a field \(k\). Let \(T \in \text{End}_k(V)\). Show that \(\text{tr}(T \otimes T) = (\text{tr}(T))^2\). Here \(\text{tr}(T)\) is the trace of \(T\).
\end{prob}
\begin{proof}
    We will show that $\tr(T\otimes T)=(\tr T)^2$, and the $T\otimes T\otimes$ is done similarly. We will use matrix representation to do an explicit computation. Let $\{v_1,\dots, v_n\}$ be a basis of$V$, then $V\otimes V$ has basis 
        \begin{equation*}
            \{v_i\otimes v_j: 1\leq i,j\leq n\}
        \end{equation*}
        and 
        \begin{equation*}
            T\otimes T(v_i\otimes v_j)=Tv_i\otimes Tv_j
        \end{equation*}
        Let $T=(a_{ij})$, then we know 
        \begin{equation*}
            (\tr(T))^2=\left(\sum_{i=1}^na_{ii}\right)^2
        \end{equation*}
        And we have 
        \begin{align*}
            T\otimes T(v_i\otimes v_j)&=\sum_{k=1}^n\sum_{l=1}^na_{ki}a_{lj}v_k\otimes v_l
        \end{align*}
        Therefore computing the trace we see 
        \begin{equation*}
            \tr(T\otimes T)=\sum_{i=1}^n\sum_{j=1}^na_{ii}a_{jj}=\tr(T)^2
        \end{equation*}
        as desired!
\end{proof}

\begin{prob}[S2016-Q4]
    Let \(V\) and \(W\) be two finite dimensional vector spaces over a field \(K\). Show that for any \(q>0\),
    \[\bigwedge^{q}(V\oplus W)\cong\sum_{i=0}^{q}(\bigwedge^{i}(V)\otimes_{K}\bigwedge^{q-i}(W)).\]
\end{prob}
\begin{proof}
    Any two finite dimensional vector spaces of the same dimension are isomorphic. Hence, it suffices to show that the dimensions are equal. We will convince ourselves it holds for $q=2$. Let $\{v_1,\dots, v_n\}$ be the basis of $V$, and $\{w_1,\dots,w_k\}$ be the basis of $W$, then we begin with the LHS:
    \begin{equation*}
        \bigwedge^{2}(V\oplus W)
    \end{equation*}
    We note that $V\oplus W$ has basis 
    \begin{equation*}
        \left\{(v_i,w_j):1\leq i\leq n, 1\leq j\leq k\right\}
    \end{equation*}
    So we reenumerate the $n+k$ basis as 
    \begin{equation*}
        \{e_1,\dots, e_{n+k}\}
    \end{equation*}
    Then $\bigwedge^{q}(V\oplus W)$ has basis
    \begin{equation*}
        \left\{ e_i\wedge e_j: i<j\right\}
    \end{equation*}
    There are exactly $1+\dots+(n+k-1)$ basis vectors i.e.,
    \begin{equation*}
        \dim\left(\bigwedge^{2}(V\oplus W)\right)=\frac{(n+k-1)(n+k)}{2}
    \end{equation*}
    As for the RHS: 
    \begin{equation*}
        \dim\left( \sum_{i=0}^{2}(\bigwedge^{i}(V)\otimes_{K}\bigwedge^{2-i}(W))\right)
        \frac{(k-1)k}{2}+nk+\frac{(n-1)n}{2}
    \end{equation*}
    And we observe that two two quantities are equal. Now we do the general case, just like above, 
    \begin{equation*}
        \dim\left(\bigwedge^{q}(V\oplus W)\right)=\binom{n+k}{q}
    \end{equation*}
    And the RHS:
    \begin{equation*}
        \dim\left(\bigwedge^{q-1}(V\oplus W)\wedge (V\oplus W)\right)=\sum_{i=0}^q\binom{n}{i}\binom{k}{q-i}
    \end{equation*}
    and it is clear that these two quantities are equal.
\end{proof}



\begin{prob}[S2011-Q4]
    Let \(F\) be a field, and \(V\) a finite-dimensional vector space over \(F\), with \(\dim_F V=n\).
    \begin{itemize}
        \item[(a)] Prove that if \(n>2\), the spaces \(\bigwedge^2(\bigwedge^2(V))\) and \(\bigwedge^4(V)\) are not isomorphic.
        \item[(b)] Let \(k\) be a positive integer. Prove that when \(v\in\bigwedge^k(V)\) and \(0\neq x\in V\), \(v\wedge x=0\) holds if and only if \(v=x\wedge y\) for some \(y\in\bigwedge^{k-1}(V)\).
    \end{itemize}
\end{prob}
\begin{proof}
    \begin{itemize}
        \item[(a)] This is by a dimension argument:
        \begin{equation*}
            \dim\left(\bigwedge^2(\bigwedge^2(V))\right)=\binom{\binom{n}{2}}{2}=\frac{n(n-1)(n-2)(n+1)}{2}
        \end{equation*}
        whereas 
        \begin{equation*}
            \dim\left(\bigwedge^4(V)\right)=\binom{n}{4}=\frac{n(n-1)(n-2)(n-3)}{4}
        \end{equation*}
        Thus not equal if $n>2$.
        \item[(b)] If there exists such $y$ where $v=x\wedge y$, then 
        \begin{equation*}
            v\wedge x=(x\wedge y)\wedge x=(-y\wedge x)\wedge x=0
        \end{equation*}
        Conversely, if $v=0$, then it is immediate that $v=x\wedge x$. It suffices to assume that $v\neq 0$, thus if we write 
        \begin{equation*}
            v=v_1\wedge\dots\wedge v_k
        \end{equation*}
        where $v_i$'s are distinct. Then 
        \begin{equation*}
            v\wedge x=0
        \end{equation*}
        If $v_i=\pm x$ for any $i$, we are done. If not, then we derive a contradiction: $v_1\neq x$, thus 
        \begin{equation*}
            v_1\wedge\left(v_2\wedge\dots\wedge v_k\wedge x\right)=0
        \end{equation*}
        i.e., $v_2\wedge\dots\wedge v_k\wedge x=0$, now $v_2\neq x$, and we keep going, eventually $v_k\wedge x=0$ which implies $x=\pm v_k$. 
    \end{itemize}
\end{proof}


\begin{prob}[S2010-Q4]
    Let \(V\) be a \(n\)-dimensional vector space over a field \(k\). Let \(T \in \text{End}_k(V)\).
    \begin{itemize}
        \item[(a)] Show that \(tr(T \otimes T \otimes T) = (tr(T))^3\). Here \(tr(T)\) is the trace of \(T\).
        \item[(b)] Find a similar formula for the determinant \(\det(T \otimes T \otimes T)\).
    \end{itemize}
\end{prob}
\begin{proof}
    \begin{itemize}
        \item[(a)] The trace computation is exactly the same as the one above. 
        \item[(b)] We can compute via some combinatorics:
        \begin{equation*}
            \det(T\otimes T)=(\det T)^{2n}, \det(T\otimes T\otimes T)=(\det T)^{3n^2}
        \end{equation*}
    \end{itemize}
\end{proof}








\chapter{Linear Algebra III}
Topics: random linear algebra problems

% Page 26-28

\begin{prop}
    Let $V$ be a $m$ dimensional vector space, and $W$ be $n$ dimensional. Show that $A:V\to V$ and $B: W\to W$ has 
    \begin{equation*}
        \text{Tr}(A\otimes B)=\text{Tr}(A)\text{Tr}(B)
    \end{equation*}
\end{prop}
\begin{proof}
    Use matrix representations.
\end{proof}


\begin{prob}[S2013-Q5]
    Let \(A\) and \(B\) be \(n \times n\) matrices with complex coefficients. Assume that \((A - I)^n = 0\) and \(A^k B = BA^k\) for some natural number \(k\). Prove that \(AB = BA\) (\textit{Hint}: Prove that \(A\) can be expressed as a function of \(A^k\)).
\end{prob}
\begin{proof}
    
\end{proof}

\begin{prob}[F2011-Q2]
    Consider the special orthogonal group \(G=SO(3,\mathbb{R})\), namely,
    \[G=\{A\in GL(3,\mathbb{R}): A^T A=\mathrm{I}_3, \det(A)=1\}\]
    \begin{itemize}
        \item[(a)] Show that for any element \(A\) in \(G\), there exists a real number \(\alpha\) with \(-1\leq\alpha\leq 3\) such that
        \[A^3-\alpha A^2+\alpha A-\mathrm{I}_3=0.\]
        \item[(b)] For which real numbers \(\alpha\) with \(-1\leq\alpha\leq 3\) does there exist an element \(A\) in \(G\) whose minimal polynomial is \(x^3-\alpha x^2+\alpha x-1\)? Explain your answer.
    \end{itemize}
\end{prob}
\begin{proof}
    \begin{itemize}
        \item[(a)] The determinant forces the eigenvalues (over $\C$) to have norm $1$. The form is done by explicit computations.
        \item[(b)] It has the minimal polynomial equal to the characteristic polynomial if the polynomial splits into three distinct roots, we know $x=1$ has a root, 
        \begin{equation*}
            (x-1)(x^2+(1-\alpha)x+1)
        \end{equation*}
        Hence as long as $\alpha\neq -1,3$, the minimal polynomial and the characteristic polynomial coincide.
    \end{itemize}
\end{proof}


\begin{prob}[F2007-Q3]
    Let \(A = \begin{pmatrix} a & b \\ c & d \end{pmatrix}\) be a real matrix such that \(a,b,c,d > 0\).
    \begin{itemize}
        \item[(1)] Prove that \(A\) has two distinct real eigenvalues, \(\lambda > \mu\).
        \item[(2)] Prove that \(\lambda\) has an eigenvector in the first quadrant and \(\mu\) has an eigenvector in the second quadrant.
    \end{itemize}
\end{prob}






\begin{prob}[S2007-Q1]
    Prove that the integer orthogonal group \(O_n(\mathbb{Z})\) is a finite group. (By definition, an \(n \times n\) square matrix \(X\) over \(\mathbb{Z}\) is orthogonal if \(XX^t = I_n\).)
\end{prob}

\begin{prob}[F2008-Q4]
    A differentiation of a ring R is a mapping \(D:R\to R\) such that, for all \(x,y\in R\),
    \begin{itemize}
        \item[(1)] \(D(x+y) = D(x) + D(y)\); and
        \item[(2)] \(D(xy) = D(x)y + xD(y)\).
    \end{itemize}
    If \(K\) is a field and \(R\) is a \(K\)-algebra, then its differentiation are supposed to be over K, that is,
    \begin{itemize}
        \item[(3)] \(D(x) = 0\) for any \(x \in K\).
    \end{itemize}
    Let D be a differentiation of the K-algebra \(M_n(K)\) of \(n \times n\)-matrices. Prove that there exists a matrix \(A \in M_n(K)\) such that \(D(X) = AX - XA\) for all \(X \in M_n(K)\).
\end{prob}


\begin{prob}[F2006-Q1]
    Let \(\text{SL}_n(k)\) be the special linear group over a field \(k\), i.e, \(n \times n\) matrices with determinant 1. Let \(I\) be the identity matrix, and \(E_{ij}\) be the elementary matrix that has 1 at \((i,j)\)-entry and 0 elsewhere. Here \(1 \leq i \neq j \leq n\).
    \begin{itemize}
        \item[(1)] Let \(C_{ij}\) be the centralizer of the matrix \(I + E_{ij}\). Find explicit generators of \(C_{ij}\).
        \item[(2)] Find the intersection
        \[\bigcap_{1 \leq i \neq j \leq n} C_{ij}.\]
        \item[(3)] Determine all the elements in the conjugacy class of \(I + E_{ij}\).
    \end{itemize}
\end{prob}


\begin{prob}[S2018-Q1]
    Let \(F\) be a field of characteristic not equal to 2. Let \(D\) be the non-commutative algebra over \(F\) generated by elements \(i,j\) that satisfy the relations
    \[i^2 = j^2 = 1, \quad ij = -ji.\]
    Define \(k = ij.\)
    \begin{itemize}
        \item[(a)] Verify that \(D\) is isomorphic to the algebra \(M_2(F)\) of \(2 \times 2\) matrices in such a way that
        \[1 \leftrightarrow \begin{pmatrix} 1 & 0 \\ 0 & 1 \end{pmatrix}, i \leftrightarrow \begin{pmatrix} 1 & 0 \\ 0 & -1 \end{pmatrix}, j \leftrightarrow \begin{pmatrix} 0 & 1 \\ 1 & 0 \end{pmatrix}, k \leftrightarrow \begin{pmatrix} 0 & 1 \\ -1 & 0 \end{pmatrix}.\]
        \item[(b)] Write \(q = x + yi + zj + uk\) for \(x,y,z,u \in F\). Verify that the norm
        \[N(q) = x^2 - y^2 - z^2 + u^2\]
        corresponds to the determinant under the isomorphism of part (a).
        \item[(c)] What does the involution \(q \mapsto \bar{q} = x - yi - zj - uk\) on \(D\) correspond to on the matrix side?
    \end{itemize}
\end{prob}


\begin{prob}[S2006-Q3]
    Let \(V\) be a \(n\)-dimensional vector space over a field \(k\), with a basis \(\{e_1, \ldots, e_n\}\). Let \(A\) be the ring of all \(n \times n\) diagonal matrices over \(k\). \(V\) is a \(A\)-module under the action:
    \[\text{diag}(\lambda_1, \ldots, \lambda_n) \cdot (a_1 e_1 + \cdots + a_n e_n) = (\lambda_1 a_1 e_1 + \cdots + \lambda_n a_n e_n).\]
    Find all \(A\)-submodules of \(V\).
\end{prob}


\begin{prob}[S2006-Q1]
    Let \(\mathbb{F}_p\) be the field with \(p\) elements, here \(p\) is prime. Let \(\text{SL}_2(\mathbb{F}_p)\) be the group of \(2 \times 2\) matrices over \(\mathbb{F}_p\) with determinant 1.
    \begin{itemize}
        \item[(1)] Find the order of \(\text{SL}_2(\mathbb{F}_p)\). Deduce that
        \[H = \left\{ \begin{pmatrix} 1 & a \\ 0 & 1 \end{pmatrix} \mid a \in \mathbb{F}_p \right\}\]
        is a Sylow-subgroup of \(\text{SL}_2(\mathbb{F}_p)\).
        \item[(2)] Determine the normalizer of \(H\) in \(\text{SL}_2(\mathbb{F}_p)\) and find its order.
    \end{itemize}
\end{prob}


\begin{prob}[S2004-Q1]
    Let \(\mathbb{F}_2\) be the finite field with 2 elements.
    \begin{itemize}
        \item[(a)] What is the order of \(\text{GL}_3(\mathbb{F}_2)\), the group of \(3 \times 3\) invertible matrices over \(\mathbb{F}_2\)?
        \item[(b)] Assuming the fact that \(\text{GL}_3(\mathbb{F}_2)\) is a simple group, find the number of elements of order 7 in \(\text{GL}_3(\mathbb{F}_2)\).
    \end{itemize}
\end{prob}

\begin{prob}[S2002-Q4]
    For a field \(K\), let \(\text{SL}_2(K)\) be the special linear group over \(K\), i.e. the group of \(2 \times 2\)-matrices over \(K\) with determinant 1, and let \(\text{PSL}_2(K)\) be the quotient of \(\text{SL}_2(K)\) by its center, i.e. the projective special linear group. Find the order of \(\text{PSL}_2(F_7)\) where \(F_7\) denotes the finite field of 7 elements.
\end{prob}


\begin{prob}[S2007-Q4]
    Find the invertible elements, the zero divisors and the nilpotent elements in the following rings:
    \begin{itemize}
        \item[(a)] \(\mathbb{Z}/p^n\mathbb{Z}\), where \(n\) is a natural number, \(p\) is a prime one.
        \item[(b)] the upper triangular matrices over a field.
    \end{itemize}
\end{prob}




\chapter{Homological Algebra}

Page 29-34



\begin{prob}[S2012-Q2]
    \phantom{text}
    \begin{itemize}
        \item[(a)] Prove that if \(M\) is an abelian group and \(n\) is a positive integer, the tensor product \(M\otimes_{\mathbb{Z}}\mathbb{Z}/n\mathbb{Z}\) can be naturally identified with \(M/nM\).
        \item[(b)] Compute the tensor product over \(\mathbb{Z}\) of \(\mathbb{Z}/n\mathbb{Z}\) with each of \(\mathbb{Z}/m\mathbb{Z}\), \(\mathbb{Q}\) and \(\mathbb{Q}/\mathbb{Z}\). Also compute the tensor products \(\mathbb{Q}\otimes_{\mathbb{Z}}\mathbb{Q}\), \(\mathbb{Q}\otimes_{\mathbb{Z}}(\mathbb{Q}/\mathbb{Z})\), and \((\mathbb{Q}/\mathbb{Z})\otimes_{\mathbb{Z}}(\mathbb{Q}/\mathbb{Z})\).
        \item[(c)] Let \(\mathbb{Z}^\mathbb{N}\) denote the (abelian) group of sequences \((a_i)_{i\in\mathbb{N}}\) in \(\mathbb{Z}\) under termwise addition, and \(\mathbb{Z}^{(\mathbb{N})}\) the subgroup of sequences for which \(a_i = 0\) for all but finitely many \(i\). Define \(\mathbb{Q}^\mathbb{N}\) and \(\mathbb{Q}^{(\mathbb{N})}\) analogously. Compare \(\mathbb{Z}^{(\mathbb{N})}\otimes_{\mathbb{Z}}\mathbb{Q}\) to \(\mathbb{Q}^{(\mathbb{N})}\), and \(\mathbb{Z}^\mathbb{N}\otimes_{\mathbb{Z}}\mathbb{Q}\) to \(\mathbb{Q}^\mathbb{N}\).
    \end{itemize}
\end{prob}

\begin{prob}[F2006-Q4]
    Let \(R\) be a commutative ring. Let \(M\) be an \(R\)-module.
    \begin{itemize}
        \item[(1)] Write down the definition of \(\mathcal{T}(M)\), the tensor algebra of \(M\).
        \item[(2)] Assume \(R = \mathbb{Z}\) and \(M = \mathbb{Q}/\mathbb{Z}\). Compute \(\mathcal{T}(M)\).
        \item[(3)] If \(M\) is a vector space over a field \(R\), show that \(\mathcal{T}(M)\) contains no zero divisors.
    \end{itemize}
\end{prob}

\begin{prob}[S2009-Q5]
    Consider the \(\mathbb{Z}\)-modules \(M_i = \mathbb{Z}/2^i\mathbb{Z}\) for all positive integers \(i\). Let \(M = \prod_{i=1}^{\infty} M_i\). Let \(S = \mathbb{Z} - \{0\}\).
    \begin{itemize}
        \item[(a)] Show that
        \[\mathbb{Q} \otimes_{\mathbb{Z}} M \cong S^{-1}M.\]
        Here \(S^{-1}M\) is the localization of \(M\).
        \item[(b)] Show that
        \[\mathbb{Q} \otimes_{\mathbb{Z}} \prod_{i=1}^{\infty} M_i \neq \prod_{i=1}^{\infty} (\mathbb{Q} \otimes_{\mathbb{Z}} M_i).\]
    \end{itemize}
\end{prob}

\begin{prob}[S2013-Q1]
    Prove that, as a $\mathbb{Z}$-module, $\mathbb{Q}$ is flat but not projective.
\end{prob}


\begin{prob}[F2008-Q5]
    For each \(n \in \mathbb{Z}\), define the ring homomorphism
    \[\phi_n : \mathbb{Z}[x] \to \mathbb{Z} \text{ by } \phi_n(f) = f(n).\]
    This gives a \(\mathbb{Z}[x]\)-module structure on \(\mathbb{Z}\), i.e,
    \[f \circ a = f(n) \cdot a \text{ for all } f \in \mathbb{Z}[x] \text{ and } a \in \mathbb{Z}.\]
    Now given two integers \(m,n \in \mathbb{Z}\), compute the tensor product \(\mathbb{Z} \otimes_{\mathbb{Z}[x]} \mathbb{Z}\) where the left-hand copy of \(\mathbb{Z}\) uses the module structure from \(\phi_n\) and the right-hand copy of \(\mathbb{Z}\) uses the module structure from \(\phi_m\). (Note: The answer depends on the numbers \(n\) and \(m\).)
\end{prob}


\begin{prob}[F2014-Q2]
    Let \(R = \mathbb{Q}[X]\), \(I\) and \(J\) the principal ideals generated by \(X^2 - 1\) and \(X^3 - 1\) respectively. Let \(M = R/I\) and \(N = R/J\). Express in simplest terms [the isomorphism type of] the \(R\)-modules \(M \otimes_R N\) and \(\text{Hom}_R(M, N)\). \textbf{Explain.}
\end{prob}


\begin{prob}[F2004-Q5]
    Consider the ideal \(I = (2,x)\) in \(R = \mathbb{Z}[x]\).
    \begin{itemize}
        \item[(a)] Construct a non-trivial \(R\)-module homomorphism \(I \otimes_R I \to R/I\), and use that to show that \(2 \otimes x - x \otimes 2\) is a non-zero element in \(I \otimes_R I\).
        \item[(b)] Determine the annihilator of \(2 \otimes x - x \otimes 2\).
    \end{itemize}
\end{prob}


\begin{prob}[S2018-Q5]
    Let \(n\) be a positive integer and \(A\) an abelian group. Prove that
    \[\text{Ext}^1(\mathbb{Z}/n\mathbb{Z}, A) \cong A/nA.\]
\end{prob}


\begin{prob}[F2002-Q3]
    Working over the integers, calculate (and show your work in a readable fashion) \(\text{Tor}(\mathbb{Z}/(p), \mathbb{Z}/(p))\).
\end{prob}

\begin{prob}[F2002-Q4]
    Working over the integers, calculate (and show your work in a readable fashion) \(\text{Ext}(\mathbb{Z}/(p), \mathbb{Z}/(p))\).
\end{prob}



\begin{prob}[S2018-Q2]
    Let \(R\) be a commutative ring. An \(R\)-module \(M\) is said to be finitely presented if there exists a right-exact sequence
    \[R^m \longrightarrow R^n \longrightarrow M \longrightarrow 0\]
    for some non-negative integers \(m,n\). Prove that any finitely generated projective \(R\)-module \(P\) is finitely presented.
\end{prob}



\begin{prob}[F2013-Q3]
    Let \(R\) be a commutative ring with unity. Given an \(R\)-module \(A\) and an ideal \(I \subset R\), there is a natural \(R\)-module homomorphism \(A \otimes_R I \to A \otimes_R R \cong A\) induced by the inclusion \(I \subset R\). In the following three steps you shall prove the flatness criterion: \textit{A is flat if and only if for every finitely generated ideal \(I \subset R\) the natural map \(A \otimes_R I \to A \otimes_R R\) is injective.}
    \begin{itemize}
        \item[(a)] Prove that if \(A\) is flat and \(I \subset R\) is a finitely generated ideal then \(A \otimes_R I \to A \otimes_R R\) is injective.
        \item[(b)] If \(A \otimes_R I \to A \otimes_R R\) is injective for every finitely generated ideal \(I\), prove that \(A \otimes_R I \to A \otimes_R R\) is injective for every ideal \(I\). Show that if \(K\) is any submodule of a free module \(F\) then the natural map \(A \otimes_R K \to A \otimes_R F \cong A\) induced by the inclusion \(K \subset F\) is injective (\textit{Hint}: the general case reduces to the case when \(F\) has finite rank).
        \item[(c)] Let \(\psi : L \to M\) be an injective homomorphism of \(R\)-modules. Prove that the induced map \(1 \otimes \psi : A \otimes_R L \to A \otimes_R M\) is injective (\textit{Hint}: Write \(M\) as a quotient \(f : F \to M\) of a free module \(F\), giving a short exact sequence \(0 \to K \to F \to M \to 0\) and consider the commutative diagram
        \[\begin{tikzcd}
            0 & K & J & L & 0 \\
            0 & K & F & M & 0
            \arrow[from=1-1, to=1-2]
            \arrow[from=1-2, to=1-3]
            \arrow["{\text{id}}", from=1-2, to=2-2]
            \arrow[from=1-3, to=1-4]
            \arrow[from=1-4, to=1-5]
            \arrow["\varphi", from=1-4, to=2-4]
            \arrow[from=2-1, to=2-2]
            \arrow[from=2-2, to=2-3]
            \arrow["f", from=2-3, to=2-4]
            \arrow[from=2-4, to=2-5]
        \end{tikzcd}\]
        where \(J = f^{-1}(\psi(L))\)).
    \end{itemize}
\end{prob}



\begin{prob}[F2013-Q4]
    \phantom{text}
    \begin{itemize}
        \item[(a)] Let \(R\) be a P.I.D. Prove that a finitely generated \(R\)-module \(M\) is flat if and only if \(M\) is torsion-free (hence, free by the structure theorem).
        \item[(b)] Give an example of an integral domain \(R\) and a torsion-free \(R\)-module \(M\) such that \(M\) is not free.
    \end{itemize}
\end{prob}

\begin{prob}[F2000-Q6]
    Let \(R\) be the ring \(\mathbb{Q}[X]/(X^7-1)\), where \((X^7-1)\) is the ideal generated by \(X^7-1\) in \(\mathbb{Q}[X]\). Give an example of a finitely generated projective \(R\)-module which is not \(R\)-free. (We remind you that an \(R\)-module is called projective if it is a direct summand of a free \(R\)-module.)
\end{prob}










\chapter{Commutative Algebra}
Topics: basic properties, Nakayama's lemma, integrality.

Page 35-40


\begin{prob}[S2017-Q1]
    Let \(A\) be a commutative ring, and define the \textit{nilradical} \(\sqrt{0}\) to be the set of nilpotent elements in \(A\). Show that \(\sqrt{0}\) is equal to the intersection of all prime ideals in \(A\). Show that if \(A\) is reduced, then \(A\) can be embedded into a product of fields.
\end{prob}

\begin{prob}[F2004-Q2]
    Let \(\mathfrak{N}\) be the set of all nilpotent elements in a ring \(R\). Assume first that \(R\) is commutative.
    \begin{itemize}
        \item[(a)] Show that \(\mathfrak{N}\) is an ideal in \(R\), and \(R/\mathfrak{N}\) contains no non-zero nilpotent elements.
        \item[(b)] Show that \(\mathfrak{N}\) is the intersection of all the prime ideals of \(R\).
        \item[(c)] Give an example with \(R\) \textbf{non}-commutative where \(\mathfrak{N}\) is not an ideal in \(R\).
    \end{itemize}
\end{prob}

\begin{prob}[S2014-Q4]
    Let \(L/K\) be a Galois extension of degree \(p\) with \(\text{char}K=p\). Show that \(L=K(\theta)\), where \(\theta\) is a root of \(x^{p}-x-a,a\in K\), and, conversely, any such extension is Galois of degree 1 or \(p\).
\end{prob}





\begin{prob}[S2009-Q2]
    Consider \(\mathbb{Z}[\omega] = \{a + b\omega \mid a, b \in \mathbb{Z}\}\) where \(\omega\) is a non-trivial cube root of 1. Show that \(\mathbb{Z}[\omega]\) is an Euclidean domain.
\end{prob}


\begin{prob}[F2006-Q3]
    Let \(A\) be a principal integral domain and \(K\) be its field of fractions. Assume that \(R\) is a ring such that \(A \subset R \subset K\). Show that \(R\) is also a principal integral domain.
\end{prob}


\begin{prob}[F2001-Q2]
    Let \(S\) denote the ring \(\mathbb{Z}[X]/X^2\mathbb{Z}[X]\), where \(X\) is a variable.
    \begin{itemize}
        \item[(a)] Show that \(S\) is a free \(\mathbb{Z}\)-module and find a \(\mathbb{Z}\)-basis for \(S\).
        \item[(b)] Which elements of \(S\) are units (i.e. invertible with respect to multiplication)?
        \item[(c)] List all the ideals of \(S\).
        \item[(d)] Find all the nontrivial ring morphisms defined on \(S\) and taking values in the ring of Gaussian integers \(\mathbb{Z}[i]\).
    \end{itemize}
\end{prob}

\begin{prob}[S2001-Q6]
    Let \(R\) be the ring \(\mathbb{Z}[X,Y]/(YX^2-Y)\), where \(X\) and \(Y\) are two algebraically independent variables, and \((YX^2-Y)\) is the ideal generated by \(YX^2-Y\) in \(\mathbb{Z}[X,Y]\).
    \begin{itemize}
        \item[(a)] Show that the ideal \(I\) generated by \(Y-4\) in \(R\) is not prime.
        \item[(b)] Provide the complete list of prime ideals in \(R\) containing the ideal \(I\) described in question (a).
        \item[(c)] Which of the ideals found in (b) are maximal?
    \end{itemize}
\end{prob}

\begin{prob}[F2017-Q3]
    In this problem all rings are commutative.
    \begin{itemize}
        \item[(a)] Let \(R\) be a local ring with maximal ideal \(\mathfrak{m}\), let \(N\) and \(M\) be finitely generated \(R\)-modules, and let \(f\colon N\to M\) be an \(R\)-linear map. Show that \(f\) is surjective if and only if the induced map \(N/\mathfrak{m}N\to M/\mathfrak{m}M\) is.
        \item[(b)] Recall that a module \(M\) over a ring \(R\) is \textit{projective} if the functor \(\operatorname{Hom}_{R}(M,-)\) is exact. Show that if \(R\) is local and \(M\) is finitely generated projective, then \(M\) is free.
    \end{itemize}
\end{prob}

\begin{prob}[F2010-Q4]
    Let \(A\) be a commutative Noetherian local ring with maximal ideal \(\mathfrak{m}\). Assume \(\mathfrak{m}^n = \mathfrak{m}^{n+1}\) for some \(n > 0\). Show that \(A\) is Artinian.
\end{prob}

\begin{prob}[F2009-Q5]
    Let \(A, B\) be two Noetherian local rings with maxima ideals \(m_A, m_B\), respectively. Let \(f : A \to B\) be a ring homomorphism such that \(f^{-1}(m_B) = m_A\). Assume that:
    \begin{itemize}
        \item[1.] \(A/m_A \to B/m_B\) is an isomorphism.
        \item[2.] \(m_A \to m_B/m_B^2\) is surjective.
        \item[3.] \(B\) is a finitely generated \(A\)-module (via \(f\)). Show that \(f\) is surjective.
    \end{itemize}
\end{prob}


\begin{prob}[F2015-Q6]
    Let \(K\) be a finite algebraic extension of \(\mathbb{Q}\).
    \begin{itemize}
        \item[(a)] Give the definition of an integral element of \(K\).
        \item[(b)] Show that the set of integral elements in \(K\) form a sub-ring of \(K\).
        \item[(c)] Determine the ring of integers in each of the following two fields No credit for memorized answers: \(\mathbb{Q}(\sqrt{13})\), and \(\mathbb{Q}(\sqrt[3]{2})\).
    \end{itemize}
\end{prob}

\begin{prob}[F2009-Q2]
    Consider \(\mathbb{Q}[\sqrt{5}] = \{a + b\sqrt{5}|a,b \in \mathbb{Q}\}\). Determine the integral closure of \(\mathbb{Z}\) in \(\mathbb{Q}[\sqrt{5}]\).
\end{prob}

\begin{prob}[S2012-Q5]
    \phantom{text}
    \begin{itemize}
        \item[(a)] Give the definition of a Dedekind domain.
        \item[(b)] Give an example of a Dedekind domain that is not a principal ideal domain. Verify from the definition that it \textit{is} a Dedekind domain, and also that it isn't a principal ideal domain.
    \end{itemize}
\end{prob}


\begin{prob}[S2005-Q5]
    Let \(A\) be an integral domain and let \(K\) be its field of fractions. Let \(A'\) be the integral closure of \(A\) in \(K\). Let \(P \subset A\) be a prime ideal and let \(S = A - P\). (Note that \(A_P = S^{-1}A\) is contained in \(K\).) Show that \(A_P\) is integrally closed in \(K\) if and only if \((A'A) \otimes_A A_P = 0\).
\end{prob}


\begin{prob}[F2013-Q2]
    Let \(a\) be an integral algebraic number such that its norm is 1 for any imbedding into \(\mathbb{C}\), the field of complex numbers. Prove that \(a\) is a root of unity.
\end{prob}

\begin{prob}[F2004-Q4]
    Let \(\lambda_1, \ldots, \lambda_n\) be roots of unity, with \(n \geq 2\). Assume that \(\frac{1}{n} \sum_{i=1}^n \lambda_i\) is integral over \(\mathbb{Z}\). Show that either \(\sum_{i=1}^n \lambda_i = 0\) or \(\lambda_1 = \lambda_2 = \cdots = \lambda_n\).
\end{prob}














\chapter{Ring Theory Random}
Page 41

\begin{prob}[S2010-Q2]
    Let \(R\) be a ring such that \(r^3 = r\) for all \(r \in R\). Show that \(R\) is commutative. (Hint: First show that \(r^2\) is central for all \(r \in R\).)
\end{prob}

\begin{prob}[S2006-Q2]
    Let \(R\) be a ring with identity 1. Let \(x, y \in R\) such that \(xy = 1\).
    \begin{itemize}
        \item[(1)] Assume \(R\) has no zero-divisor. Show that \(yx = 1\).
        \item[(2)] Assume \(R\) is finite. Show that \(yx = 1\).
    \end{itemize}
\end{prob}


























\chapter{Tensor Products over Fields}
Page 42-43


\begin{prob}[S2018-Q3]
    Let \(K/k\) be a finite separable field extension, and let \(L/k\) be any field extension. Show that \(K\otimes_{k}L\) is a product of fields.
\end{prob}

\begin{prob}[F2019-Q3]
    Let \(F,L\) be extensions of a field \(K\). Suppose that \(F/K\) is finite. Show that there exists an extension \(E/K\) such that there are monomorphisms of \(F\) into \(E\) and of \(L\) into \(E\) which are identical on \(K\).
\end{prob}


\begin{prob}[F2009-Q4]
    Let \(E\) and \(F\) be finite field extensions of a field \(k\) such that \(E \cap F = k\), and that \(E\) and \(F\) are both contained in a larger field \(L\). Assume that \(E\) is Galois over \(k\). Show that \(E \otimes_k F \cong EF\).
\end{prob}



\begin{prob}[S2008-Q5]
    Let \(k\) be a field of characteristic zero. Assume that \(E\) and \(F\) are algebraic extensions of \(k\) and both contained in a larger field \(L\). Show that the \(k\)-algebra \(E \otimes_k F\) has no nonzero nilpotent elements.
\end{prob}


\begin{prob}[S2004-Q5]
    Show that there is a \(\mathbb{C}\)-algebra isomorphism between \(\mathbb{C} \otimes_{\mathbb{R}} \mathbb{C}\) and \(\mathbb{C} \times \mathbb{C}\).
\end{prob}


\begin{prob}[F2005-Q5]
    Let \(\mathbb{C}\) and \(\mathbb{R}\) be complex and real number fields. Let \(\mathbb{C}(x)\) and \(\mathbb{C}(y)\) be function fields of one variable. Consider \(\mathbb{C}(x) \otimes_{\mathbb{R}} \mathbb{C}(y)\) and \(\mathbb{C}(x) \otimes_{\mathbb{C}} \mathbb{C}(y)\).
    \begin{itemize}
        \item[(1)] Determine if they are integral domains.
        \item[(2)] Determine if they are fields.
    \end{itemize}
\end{prob}

\begin{prob}[F2003-Q4]
    Verify the isomorphism of algebras over a field \(K\):
    \[\mathbb{M}_n(K) \otimes_K \mathbb{M}_{m}(K) \simeq \mathbb{M}_{mn}(K).\]
    [Note: \(\mathbb{M}_n(K)\) denotes the algebra of \(n \times n\) matrices over \(K\).]
\end{prob}

















\chapter{Irreducibility of Polynomials}
Page 44-45

\begin{prob}[F2008-Q2]
    Show that the polynomial \(x^5 - 5x^4 - 6x - 2\) is irreducible in \(\mathbb{Q}[x]\).
\end{prob}

\begin{prob}[F2003-Q3]
    Obtain a factorization into irreducible factors in \(\mathbb{Z}[x]\) of the polynomial \(x^{10} - 1\).
\end{prob}


\begin{prob}[S2004-Q3]
    Let \(k\) be a field with characteristic 0. Let \(m \geq 2\) be an integer. Show that \(f(x,y) = x^m + y^m + 1\) is irreducible in \(k[x,y]\).
\end{prob}

\begin{prob}[S2017-Q2, S2007-Q3]
    Write down the minimal polynomial for \(\sqrt{2}+\sqrt{3}\) over \(\mathbb{Q}\) and prove that it is reducible over \(\mathbb{F}_{p}\) for every prime number \(p\).
\end{prob}

\begin{prob}[S2015-Q4]
    Prove that the polynomial \(x^{4}+1\) is not irreducible over any field of positive characteristic.
\end{prob}



\begin{prob}[F2010-Q2]
    \phantom{text}
    \begin{itemize}
        \item[(a)] Find the complete factorization of the polynomial \(f(x) = x^6 - 17x^4 + 80x^2 - 100\) in \(\mathbb{Z}[x]\).
        \item[(b)] For which prime numbers \(p\) does \(f(x)\) have a root in \(\mathbb{Z}/p\mathbb{Z}\) (i.e, \(f(x)\) has a root modulo \(p\))? Explain your answer.
    \end{itemize}
\end{prob}

















\chapter{Galois Theory}
Page 46
Page 52-61



Quick reminder whether a polynomial has a rational root:
\begin{prop}
    Let $f(t)=a_nt^n+\dots+a_1t+a_0$, and if a rational (expressed in lowest terms) $\frac{p}{q}$ is a root of $f$, then $p\mid a_0, q\mid a_0$.
\end{prop}


\begin{defn}[Galois extension]
    A field extension $k\subset L$ is Galois if for all $x\in L$, the minimal polynomial $f(x)\in k[x]$ splits into a linear factor without repeated roots.
\end{defn}
\begin{defn}[normal extension]
    An extension $k\subset K$ is normal if $f$ has a root in $K$ if and only if $f$ splits completely into linear factors over $K$. An extension that is normal and separable is Galois.
\end{defn}

\begin{thm}
    Suppose $k\subset L$ is Galois, 
    \begin{equation*}
        \left\{ k\subset M\subset L\right\} \xLeftrightarrow{\text{ one-to-one}} \left\{\text{Subgroups of Gal}(L/k) \right\}
    \end{equation*}
    Moreover, the order of the Galois group is the degree of the field extension.
    \begin{equation*}
        \left|\text{Gal}(L/k)\right|=[L:k]
    \end{equation*}
\end{thm}

% \begin{example}
%     Let $f(x)=x^4-4x^2-2\in\Q[x]$, compute the Galois group of $f$.
% \end{example}


\begin{prop}
    Let $G$ be a Galois group of a polynomial $f$ of degree $4$, and $|G|=8$, then 
    \begin{equation*}
        G\cong D_8
    \end{equation*}
\end{prop}
\begin{proof}
    We know that $G$ permutes the four roots of $f$, i.e., $G$ embeds into $S_4$. Since $|G|=8$, we know $G$ is a Sylow-$2$ subgroup of $S_4$, and all Sylow-$2$ subgroups are conjugates (isomorphic to one another), i.e., 
    \begin{equation*}
        G\cong D_8
    \end{equation*}
    as desired.
\end{proof}

\begin{prop}
    Let $k\subset K$ be a Galois extension, then the intermediate field extensions $k\subset E\subset K$ is determined by the subgroups of $\gal(K/k)$. Namely, let $E$ be an intermediate extension, there exists a subgroup $H$ of $\gal(K/k)$ that fixes $E$. This extension is normal if and only if $H$ is normal. And $E/k$ is Galois if and only if $H$ is normal.
\end{prop}

\begin{prop}
    We can draw the lattice of subgroups of $\gal(K/k)$ and lattice of subfields:

    \begin{tikzpicture}[node distance=2cm]
        % Nodes for fields
        \node (k) {$k$};
        \node[right of=k] (E) {$E$};
        \node[right of=E] (F) {$F$};
        
        % Nodes for automorphism groups (placed below)
        \node[below of=k, yshift=-1.5cm] (Autk) {$\text{Aut}_k(F)$};
        \node[below of=E, yshift=-1.5cm] (AutE) {$\text{Aut}_E(F)$};
        \node[below of=F, yshift=-1.5cm] (e) {$\{e\}$};
        
        % Field extension lines
        \draw (k) -- (E) node[midway, above] {$\subseteq$};
        \draw (E) -- (F) node[midway, above] {$\subseteq$};
        
        % Automorphism group lines (reverse inclusions)
        \draw (Autk) -- (AutE) node[midway, below] {$\supseteq$};
        \draw (AutE) -- (e) node[midway, below] {$\supseteq$};
        
        % Vertical dotted lines to show correspondence
        \draw[dashed] (k) -- (Autk);
        \draw[dashed] (E) -- (AutE);
        \draw[dashed] (F) -- (e);
    \end{tikzpicture}
    

\end{prop}




\begin{prob}[S2009-Q3]
    Consider the field \(K = \mathbb{Q}(\sqrt{a})\) where \(a \in \mathbb{Z}, a < 0\). Show that \(K\) cannot be embedded in a cyclic extension whose degree over \(\mathbb{Q}\) is divisible by 4.
\end{prob}
\begin{proof}
    Suppose $K$ embedes into a degree $4n$ extension $L$, and 
    \begin{equation*}
        \text{Gal}(L/\Q)=\frac{\Z}{4n\Z}
    \end{equation*} 
    Since $K$ is a degree $2$ extension of $\Q$, thus $L/K$ is a degree $4n/2$ Galois extension, with Galois group 
    \begin{equation*}
        \text{Gal}(L/K)=\frac{2\Z}{4n\Z}
    \end{equation*}
    We notice that $\sqrt{a}$ is complex, hence the complex conjugation $\tau$ is in $\text{Gal}(L/\Q)$, i.e., it is an order $2$ element in $\frac{\Z}{4n\Z}$, it is therefore $[2n]$ i.e., 
    \begin{equation*}
        \tau\in\frac{2\Z}{4n\Z}=\text{Gal}(L/\Q)
    \end{equation*}
    This implies $\tau$ fixed $K$, however $\tau(\sqrt{a})\neq\sqrt{a}$, hence a contradiction.
\end{proof}


\begin{prob}[F2000-Q4]
    Let \(G\) be a finite group. Show that there exists a Galois field extension \(K/k\) whose Galois group is isomorphic to \(G\).
\end{prob}


\begin{prob}[S2001-Q2]
    Let \(K\) be the splitting field of \(f(X)=X^3-2\) over \(\mathbb{Q}\).
    \begin{itemize}
        \item[(a)] Determine an explicit set of generators for \(K\) over \(\mathbb{Q}\).
        \item[(b)] Show that the Galois group \(G(K/\mathbb{Q})\) of \(K\) over \(\mathbb{Q}\) is isomorphic to the symmetric group \(S_3\).
        \item[(c)] Provide the complete list of intermediate fields \(k\), \(\mathbb{Q} \subseteq k \subseteq K\), satisfying \([k:\mathbb{Q}]=3\).
        \item[(d)] Which of the fields determined in (c) are normal extensions of \(\mathbb{Q}\)?
    \end{itemize}
\end{prob}
\begin{proof}
    \begin{enumerate}
        \item[(a)] The set of generators is 
        \begin{equation*}
            \left\{\sqrt[3]{2}, e^\frac{2\pi i}{3} \right\}
        \end{equation*}
        \item[(b)] The minimal polynomial of $e^\frac{2\pi i}{3}$ is $x^2+x+1$. This shows that the Galois group $G$ has order $6$, and because of the complex root, there exists an element of order $2$, a transposition, that only swaps the two complex roots, and $G$ also has an element of order $3$ because $3$ divides $|G|$, this shows that $G$ must be $S_3$.
        \item[(c)]  The following is the subgroup lattice of $S_3$ and subfield lattice: 
        \[\begin{tikzcd}
            & K &&&& {\{e\}} \\
            {\mathbb{Q}(\sqrt[3]{2})} & {\mathbb{Q}(\sqrt[3]{2}\omega)} & {\mathbb{Q}(\sqrt[3]{2}\omega^2)} & {\mathbb{Q}(\omega)} & {\langle(23)\rangle} & {\langle(13)\rangle} & {\langle(12)\rangle} & {\langle(123)\rangle} \\
            & {\mathbb{Q}} &&&& {S_3}
            \arrow["2"', no head, from=1-2, to=2-3]
            \arrow["3", no head, from=1-6, to=2-7]
            \arrow["2", curve={height=-12pt}, no head, from=1-6, to=2-8]
            \arrow["2", no head, from=2-1, to=1-2]
            \arrow["3"', no head, from=2-1, to=3-2]
            \arrow["2"', no head, from=2-2, to=1-2]
            \arrow["3", no head, from=2-2, to=3-2]
            \arrow["3"', no head, from=2-4, to=1-2]
            \arrow["3", no head, from=2-5, to=1-6]
            \arrow["2"', no head, from=2-5, to=3-6]
            \arrow["3"', no head, from=2-6, to=1-6]
            \arrow["3", no head, from=3-2, to=2-3]
            \arrow["2", no head, from=3-2, to=2-4]
            \arrow["2"', no head, from=3-6, to=2-6]
            \arrow["2", no head, from=3-6, to=2-7]
            \arrow["3"', no head, from=3-6, to=2-8]
        \end{tikzcd}\]
        Thus all the $\Q\subset k$ such that $[k:\Q]=3$ are 
        \begin{equation*}
            \{\Q(\sqrt[3]{2}), \Q(\sqrt[3]{2})e^\frac{2\pi i}{3}, \Q(\sqrt[3]{2})e^\frac{4\pi i}{3}\}
        \end{equation*}


        % \\[
        %     \begin{array}{c}
        %     S_3 \\
        %     \hline
        %     \begin{array}{cccc}
        %      & A_3 & & \\
        %     \langle (1\,2) \rangle & \langle (1\,3) \rangle & \langle (2\,3) \rangle & \\
        %      & \{e\} & & \\
        %     \end{array} \\
        %     \end{array}
        %     \]
        \item[(d)] None of the above are normal because the subgroups 
        \begin{equation*}
            \{\la(12)\ra, \la(13)\ra, \la(23)\ra\}
        \end{equation*}
        are all Sylow $2$-subgroups of $S_3$, hence all conjugates to one another, i.e., not normal. 
    \end{enumerate}
\end{proof}


\begin{prob}[F2001-Q4]
    Let \(K:=\mathbb{Q}(\sqrt{3}+\sqrt{5})\).
    \begin{itemize}
        \item[(a)] Show that \(K\) is the splitting field of \(X^4-6X^2+4\).
        \item[(b)] Find the structure of the Galois group of \(K/\mathbb{Q}\).
        \item[(c)] List all the fields \(k\), satisfying \(\mathbb{Q} \subseteq k \subseteq K\).
    \end{itemize}
\end{prob}
\begin{proof}
    \begin{itemize}
        \item[(a)] I belive there is typo in (a) where the polynomial should be $f(X)=X^4-16X^2+4$. This is the minimal polynomial of $\sqrt{3}+\sqrt{5}$. We see that $\Q(\sqrt{3}+\sqrt{5})=\Q(\sqrt{3},\sqrt{5})$, hence it contains all the roots of $f$.
        \item[(b)] We let $\alpha=\sqrt{3}+\sqrt{5}$, and $\beta=\sqrt{3}-\sqrt{5}$, then we see Galois group permutes 
        \begin{equation*}
            \{\alpha,-\alpha,\beta,-\beta\}
        \end{equation*}
        and we have $\alpha\beta\in\Q$. Thus just like the above, we have 
        \begin{equation*}
            \gal(K/\Q)=\frac{\Z}{2\Z}\times\frac{\Z}{2\Z}
        \end{equation*}
        \item[(c)] We know the intermediate fields are determined by the subgroup of $\Z/2\Z\times\Z/2\Z$. 
        \[\begin{tikzcd}
            & {\{e\}} \\
            {\langle(1,0)\rangle} & {\langle(1,1)\rangle} & {\langle(0,1)\rangle} \\
            & {\Z/2\Z\times\Z/2\Z}
            \arrow["2", no head, from=2-1, to=1-2]
            \arrow["2"', no head, from=2-1, to=3-2]
            \arrow["2"', no head, from=2-2, to=1-2]
            \arrow["2", no head, from=2-2, to=3-2]
            \arrow["2"', no head, from=2-3, to=1-2]
            \arrow["2", no head, from=3-2, to=2-3]
        \end{tikzcd}\]
        and let $(1,0)$ be the element such that 
        \begin{equation*}
            (1,0)\cdot (\sqrt{3}+\sqrt{5})=\sqrt{3}-\sqrt{5}
        \end{equation*}
        then we have the corresponding lattice of subfields
        \[\begin{tikzcd}
            & {\{e\}} \\
            {\langle(1,0)\rangle} & {\langle(1,1)\rangle} & {\langle(0,1)\rangle} \\
            & {\Z/2\Z\times\Z/2\Z} \\
            & {\mathbb{Q}(\sqrt{3},\sqrt{5})} \\
            {\mathbb{Q}(\sqrt{3})} & {\mathbb{Q}(\sqrt{15})} & {\mathbb{Q}(\sqrt{5})} \\
            & {\mathbb{Q}}
            \arrow["2", no head, from=2-1, to=1-2]
            \arrow["2"', no head, from=2-1, to=3-2]
            \arrow["2"', no head, from=2-2, to=1-2]
            \arrow["2", no head, from=2-2, to=3-2]
            \arrow["2"', no head, from=2-3, to=1-2]
            \arrow["2", no head, from=3-2, to=2-3]
            \arrow["2"', no head, from=5-1, to=4-2]
            \arrow["2", no head, from=5-2, to=4-2]
            \arrow["2", no head, from=5-3, to=4-2]
            \arrow["2"', no head, from=6-2, to=5-1]
            \arrow["2", no head, from=6-2, to=5-2]
            \arrow["2", no head, from=6-2, to=5-3]
        \end{tikzcd}\]
    \end{itemize}
    So all intermediate fields are 
    \begin{equation*}
        \left\{\Q(\sqrt{3}), \Q(\sqrt{15}),\Q(\sqrt{5})\right\}
    \end{equation*}
\end{proof}


% \begin{proof}
%     \begin{itemize}
%         \item[(a)] The minimal polynomial of $\sqrt{3}+\sqrt{5}$ is exactly $x^4-6x^2+4$, and it is irreducible over $\Q$, hence $K$ 
%     \end{itemize}
% \end{proof}


\begin{prob}[F2013-Q5]
    Compute the Galois group of \(f(x) = x^4 + 1\) over \(\mathbb{Q}\).
\end{prob}
\begin{proof}
    The splitting field for $f$ is $\Q(\xi_8)$ where $\xi_8=e^\frac{2\pi i}{8}$, and the Galois group 
    \begin{equation*}
        \text{Gal}(\Q(\xi_8)/\Q)\cong\left(\Z/8\Z\right)^\times
    \end{equation*}
    thus
    \begin{equation*}
        \left(\Z/8\Z\right)^\times\cong\frac{\Z}{2\Z}\times\frac{\Z}{2\Z}
    \end{equation*}
\end{proof}



\begin{prob}[F2016-Q4]
    \phantom{text}
    \begin{itemize}
        \item[(1)] Determine the Galois group of \(x^{4}-4x^{2}-2\) over \(\mathbb{Q}\).
        \item[(2)] Let \(G\) be a group of order 8 such that \(G\) is the Galois group of a polynomial of degree 4 over \(\mathbb{Q}\). Show that \(G\) is isomorphic to the Galois group in part (1).
    \end{itemize}
\end{prob}
\begin{proof}
    \begin{itemize}
        \item[(a)] There are four roots of this polynomial 
        \begin{equation*}
            \{\alpha, -\alpha, \beta, -\beta\}
        \end{equation*}
        where 
        \begin{equation*}
            \alpha=\sqrt{2+\sqrt{6}}, \beta=\sqrt{2-\sqrt{6}}
        \end{equation*}
        Thus the Galois group embeds into $S_4$. Notice that 
        \begin{equation*}
            \alpha\beta=\sqrt{2}i
        \end{equation*}
        Thus we see the Galois extension has degree $8$:
        \[\begin{tikzcd}
            {\mathbb{Q}(\sqrt{2+\sqrt{6}},\sqrt{2}i)} \\
            {\mathbb{Q}(\sqrt{2+\sqrt{6}})} \\
            {\mathbb{Q}}
            \arrow["2", no head, from=1-1, to=2-1]
            \arrow["4", no head, from=2-1, to=3-1]
        \end{tikzcd}\]
        Notice that the Galois grop $G$ is an order $8$ subgruop of $S_4$, which implies that $G$ is a Sylow $2$ subgroup, and all Sylow $2$ subgruops are isomorphic: 
        \begin{equation*}
            G\cong D_8
        \end{equation*}

        \item[(b)] The argument is given in (a).
    \end{itemize}
\end{proof}


\begin{prob}[S2008-Q3]
    Let \(K\) be the splitting field of the polynomial \(X^4 - 6X^2 - 1\) over \(\mathbb{Q}\).
    \begin{itemize}
        \item[(a)] Compute \(\text{Gal}(K/\mathbb{Q})\).
        \item[(b)] Determine all intermediate fields that are Galois over \(\mathbb{Q}\).
    \end{itemize}
\end{prob}
\begin{proof}
    \begin{itemize}
        \item[(a)] This computation is exactly same as above, as we have the four roots 
        \begin{equation*}
            \left\{\pm\sqrt{3+\sqrt{10}}, \pm\sqrt{3-\sqrt{10}}\right\}
        \end{equation*}
        and we see that $\alpha\beta=i$, thus the Galois group $\gal(K/\Q)$ has order $8$, and embeds into $S_4$, thus 
        \begin{equation*}
            \gal(K/\Q)\cong D_8
        \end{equation*}
        \item[(b)] \textcolor{red}{not finished}
    \end{itemize}
\end{proof}





\begin{prob}[S2010-Q3]
    Compute Galois groups of the following polynomials.
    \begin{itemize}
        \item[(a)] \(x^3 + t^2x - t^3\) over \(k\), where \(k = \mathbb{C}(t)\) is the field of rational functions in one variable over complex numbers \(\mathbb{C}\).
        \item[(b)] \(x^4 - 14x^2 + 9\) over \(\mathbb{Q}\).
    \end{itemize}
\end{prob}
\begin{proof}
    \begin{itemize}
        \item[(a)] The polynomial completely factors over $\C(t)$, so the Galois group is $\{e\}$.
        \item[(b)] The roots are 
        \begin{equation*}
            \left\{\pm\sqrt{7\pm 2\sqrt{10}}\right\}
        \end{equation*}
        and $\alpha\beta\in\Q$ again, hence the Galois group is $\Z/2\Z\times\Z/2\Z$.
    \end{itemize}
\end{proof}


\begin{prob}[S2013-Q6]
    Let \(K\) be the splitting field of \(x^6 - 5\) over \(\mathbb{Q}\).
    \begin{itemize}
        \item[(a)] Prove that \(x^6 - 5\) is irreducible over \(\mathbb{Q}\).
        \item[(b)] Compute the Galois group of \(K\) over \(\mathbb{Q}\).
        \item[(c)] Describe an intermediate field \(F\) such that \(F\) is not \(\mathbb{Q}\) or \(K\) and \(F/\mathbb{Q}\) is Galois.
    \end{itemize}
\end{prob}
\begin{proof}
    \begin{itemize}
        \item[(a)] By Eisenstein.
        \item[(b)] We know $K=\Q(\sqrt[6]{5}, \zeta_6)$, where $\zeta_6$ is the $6$th root of unity. The roots are 
        \begin{equation*}
            \left\{\sqrt[6]{5}, \sqrt[6]{5}\zeta_6, \dots, \sqrt[6]{5}\zeta_6^5\right\}
        \end{equation*}
        Note that the minimal polynomial for $\zeta_6$ is $x^2-x+1$, so the size of $\gal(K/\Q)$ is $12$. We see that any $\sigma\in \gal(K/\Q)$ is determined by where it sends $\sqrt[6]{5}$ and $\zeta_6$, so we  only need to compute the possibilities of them. The Galois action is transitive implies that there $\sqrt[6]{5}$ can be sent to any $\sqrt[6]{5}\zeta_6^k$, where $k=0, 1,2,3,4,5$, and since $\zeta_6$ has minimal polynomial 
        \begin{equation*}
            x^2-x+1
        \end{equation*}
        Then there are two possibilities for $\zeta_6\mapsto \zeta_6, \bar{\zeta}_6$, where $\bar{\zeta_6}=\zeta_6^5$. Now we see that 
        \begin{equation*}
            \gal(K/Q)=D_{12}
        \end{equation*}
        as it is generated by
        \begin{equation*}
            \sigma: \sqrt[6]{5}\mapsto\zeta_6\sqrt[6]{5}, \zeta_6\mapsto\zeta_6, \quad\tau: \sqrt[6]{5}\mapsto\sqrt[6]{5},  \zeta_6\mapsto \zeta_6^5
        \end{equation*}
        satisfying $\tau\sigma=\tau\sigma^{-1}$. (One can draw a hexagon)
        \item[(c)] $F/\Q$ corresponds to a normal subgroup of $D_{12}$. Any subgroup of $6$ is normal, i.e., the subgroup 
        \begin{equation*}
            \{e, \sigma, \dots, \sigma^5\}
        \end{equation*}
        This subgroup fixes the field $\Q(\zeta_6)$. Hence it corresponds to
        \begin{equation*}
            F=\Q(\zeta_6)
        \end{equation*}
    \end{itemize}
\end{proof}





\begin{prob}[S2016-Q3]
    Determine the Galois group of \(x^{6}-10x^{3}+1\) over \(\mathbb{Q}\).
\end{prob}
\begin{proof}
    This is the same process as above, the roots are 
    \begin{equation*}
        \left\{\zeta_3^{i}\sqrt[3]{5\pm 2\sqrt{6}}: i=0,1,2\right\}
    \end{equation*}
    The order of the Galois group $G$ is $12$, but now we need another trick.
    \begin{lem}
        Transitive subgroup of $S_6$ of order 12 can only be $D_{12}$ or $A_4$. However, $A_4$ has no index $2$ subgroups, i.e., this Galois extension cannot have a subfield extension of degree $2$ over $\Q$, this gives that $G$ must be $D_{12}$:
        \[\begin{tikzcd}
            & {\mathbb{Q}(\sqrt[3]{2+\sqrt{6}},\zeta_3)} \\
            {\mathbb{Q}(\zeta_3)} & {} & {\mathbb{Q}(\sqrt[3]{2+\sqrt{6}})} \\
            & {\mathbb{Q}}
            \arrow["2"', no head, from=1-2, to=2-3]
            \arrow["6"', no head, from=2-1, to=1-2]
            \arrow["2", no head, from=2-1, to=3-2]
            \arrow["6", no head, from=3-2, to=2-3]
        \end{tikzcd}\]
    \end{lem} 
  
    % \begin{equation*}
    %     G\la \sigma,\tau: \sigma^6=\tau^2=e, \tau\sigma=\tau\sigma^{5}\ra
    % \end{equation*}
    % where 
    % \begin{equation*}
    %     \sigma: \sqrt[3]{5\pm 2\sqrt{6}}\mapsto \sqrt[3]{5\pm 2\sqrt{6}}\zeta_3, \zeta_3\mapsto\zeta_3
    % \end{equation*}
    % and 
    % \begin{equation*}
    %     \tau: \sqrt[3]{5\pm 2\sqrt{6}}\mapsto \sqrt[3]{5\pm 2\sqrt{6}}, \zeta_3\mapsto\
    % \end{equation*}
\end{proof}

\begin{prob}[F2010-Q3]
    Let \(K = \mathbb{Q}(\sqrt[5]{2}, \sqrt{-1})\) and \(F = \mathbb{Q}(\sqrt{-2})\). Show that \(K\) is Galois over \(F\) and determine the Galois group \(\text{Gal}(K/F)\).
\end{prob}
\begin{proof}

    % We know that $K$ is Galois because it is the splitting field of the following separable polynomial:
    % \begin{equation*}
    %     f(t)=t^5-2
    % \end{equation*}
    % We know that 
    % \begin{equation*}
    %     \gal(K/\Q)=S_5
    % \end{equation*}
    % because it contains a transposition and a degree $5$ element. 
\end{proof}

\begin{prob}[F2015-Q2]
    The dihedral group \(D_{2n}\) is the group on two generators \(r\) and \(s\), with respective orders \(o(r)=n\) and \(o(s)=2\), subject to the relation \(rsr=s\).
    \begin{itemize}
        \item[(a)] Calculate the order of \(D_{2n}\).
        \item[(b)] Let \(K\) be the splitting field of the polynomial \(x^8 - 2\). Determine whether the Galois group \(\text{Gal}(K/\mathbb{Q})\) is dihedral (i.e., isomorphic to \(D_{2n}\) for some \(n\)).
    \end{itemize}
\end{prob}


\begin{prob}[S2019-Q1]
    Show that any transitive subgroup of \(A_{5}\) is isomorphic to one of the following groups:
    \begin{itemize}
        \item[(a)] the cyclic group \(\mathbb{Z}/5\mathbb{Z}\),
        \item[(b)] the dihedral group \(D_{5}\),
        \item[(c)] \(A_{5}\).
    \end{itemize}
\end{prob}


\begin{prob}[S2019-Q2]
    Let \(f(x)=x^{5}-5x+12\). Verify that \(f(x)\) is irreducible in \(\mathbb{Q}[x]\) and its discriminant is \(d(f)=(2^{6}5^{3})^{2}\). If \(r_{1},\ldots,r_{5}\) are the roots of \(f\), let
    \[P(x)=\prod_{1\leq i<j\leq 5}(x-(r_{i}+r_{j})).\]
    Show that \(P(x)\) is a product of two monic irreducible polynomials in \(\mathbb{Q}[x]\):
    \[P(x)=(x^{5}-5x^{3}-10x^{2}+30x-36)(x^{5}+5x^{3}+10x^{2}+10x+4).\]
    Use this information, Problem 1 and properties of \(f_{3}\in\mathbb{F}_{3}[x]\), the reduction of \(f\) modulo 3, to show that the the Galois group \(G_{f}\) of \(f\) is isomorphic to \(D_{5}\).
\end{prob}


\begin{prob}[F2018-Q6]
    Determine the Galois group over \(\mathbb{Q}\) of the polynomial
    \[X^{6}+22X^{5}-9X^{4}+12X^{3}-37X^{2}-29X-15.\]
\end{prob}

\begin{prob}[F2017-Q4]
    Compute the Galois group of \(x^{5}-10x+5\) over \(\mathbb{Q}\).
\end{prob}


\begin{prob}[F2004-Q3]
    Let \(f(x) = x^5 - 9x + 3\). Determine the Galois group of \(f\) over \(\mathbb{Q}\).
\end{prob}

\begin{prob}[F2006-Q2]
    Let \(f\) be a polynomial in \(\mathbb{Q}[x]\). Let \(E\) be a splitting field of \(f\) over \(\mathbb{Q}\). For the following cases, determine whether \(E\) is solvable by radicals.
    \begin{itemize}
        \item[(1)] \(f(x) = x^4 - 4x + 2\).
        \item[(2)] \(f(x) = x^5 - 4x + 2\).
    \end{itemize}
\end{prob}

\begin{prob}[S2011-Q3]
    Determine the Galois group [up to isomorphism] of the splitting field of each of the following polynomials over \(\mathbb{Q}\):
    \begin{itemize}
        \item[(a)] \(f(x)=x^4-9x^3+9x+4\),
        \item[(b)] \(g(x)=x^5-6x^2+2\).
    \end{itemize}
\end{prob}


\begin{prob}[F2014-Q1]
    \phantom{text}
    \begin{itemize}
        \item[(a)] Let \(S_n\) be the symmetric group (permutation group) on \(n\) objects. Prove that if \(\sigma \in S_n\) is an \(n\)-cycle and \(\tau \in S_n\) is a transposition (i.e., a 2-cycle), then \(\sigma\) and \(\tau\) generate \(S_n\).
        \item[(b)] Let \(f_a(x)\) be the polynomial \(x^5 - 5x^3 + a\). Determine an integer \(a\) with \(-4 \leq a \leq 4\) for which \(f_a\) is irreducible over \(\mathbb{Q}\), and the Galois group of [the splitting field of] \(f_a\) over \(\mathbb{Q}\) is \(S_5\). Then explain why the equation \(f_a(x) = 0\) is not solvable in radicals.
    \end{itemize}
\end{prob}
\begin{itemize}
    \item[(a)] It suffices to assume that the $n$ cycle is $(1\dots n)$ (up to rearranging the terms), and the transposition is $(12)$. One can show that conjugation gives all the transpositions, hence generate $S_n$.
    \item[(b)] Take $a=1$, then $f_a(x)$ is irreducible: it doesn't have a root by the Rational Root Theorem and cannot be factored into lower degree polynomial by term matching. Moreover, we see that $f_a'(x)$ has $3$ roots, by Rolle's theorem, there are at most $4$ real roots, this implies that there exists a complex root $r_1$, and since this has odd degree, it must also exist a real root $r_2$. This shows that there exists an element in the Galois group that has order $5$ and a transposition (sending conjugate complex roots to each other). Thus by (a), since the Galois group is a subgroup of $S_5$, we must have it equal to $S_5$. 
\end{itemize}


\begin{prob}[F2009-Q3]
    Determine the Galois group of \(x^4 - 4x^2 + 7x - 3\) over \(\mathbb{Q}\).
\end{prob}

\begin{prob}[S2012-Q3]
    In this problem, \(G\) denotes the group \(S_5 \times C_2\), where \(S_5\) is the symmetric group on five letters and \(C_2\) is the cyclic group of order 2.
    \begin{itemize}
        \item[(a)] Determine all normal subgroups of \(G\).
        \item[(b)] Give an example of a polynomial with rational coefficients whose Galois group is \(G\), deducing that from basic principles.
    \end{itemize}
\end{prob}


\begin{prob}[F2015-Q4]
    Let \(H = S_3 \times S_5\).
    \begin{itemize}
        \item[(a)] Determine all normal subgroups of \(H\). Make sure you have them all! What would be different if \(H\) were replaced by \(S_2 \times S_5\)?
        \item[(b)] Describe, in full detail, the construction of a polynomial with rational coefficients, whose Galois group is isomorphic to \(H\).
    \end{itemize}
\end{prob}








\end{document}